\chapter{Syntactic Transition--Output Images and Algebraic Eilenberg Duality}
\label{ch:syntactic-transition-output-eilenberg}

\section{Overview and standing assumptions}
\label{sec:syntactic-to-overview-standing-assumptions}


\subsection{Orientation}
\label{subsec:syntactic-to-orientation}

The preceding chapters developed residual semantics in two related forms. Every
Mealy morphism
\[
    P:\mathbb A\rightsquigarrow\mathbb B
\]
has a unique strict semantic Mealy homomorphism
\[
    \beta_P:
    P\to\mathbf{Beh}^{\mathrm{can}}_{\mathbb A,\mathbb B},
\]
and local coequational theories are precisely the derivative-closed subobjects
\[
    V\hookrightarrow
    \mathbf{Beh}^{\mathrm{can}}_{\mathbb A,\mathbb B}.
\]
This chapter extracts from this semantic picture an algebraic
transition--output quotient of the input category.

The formal development has four stages. First, the output category
\(\mathbb B\) defines an output-comparison Kleisli monad
\[
    T_{\mathbb B}:=(t_B)_!(s_B)^*:
    \mathcal E/B_0\to\mathcal E/B_0.
\]
For a Mealy carrier
\[
    A_0\xleftarrow{p}P\xrightarrow{q}B_0,
\]
this monad yields an internal category
\[
    \operatorname{KlEnd}_{\mathbb B}(P)
\]
of fibrewise output-comparison Kleisli endomorphisms.

Second, the target-to-source convention for Mealy execution is repackaged as an
identity-on-objects internal functor
\[
    \chi_P:
    \mathbb A^{\mathrm{op}}
    \longrightarrow
    \operatorname{KlEnd}_{\mathbb B}(P).
\]
An arrow
\[
    a:u\to v
\]
of \(\mathbb A\) acts, in the opposite direction, as a Kleisli map
\[
    P_v\rightsquigarrow P_u,
    \qquad
    x\longmapsto
    \bigl(\delta(a,x),\omega(a,x)\bigr),
\]
where
\[
    \omega(a,x):q(\delta(a,x))\to q(x).
\]

Third, a local coequational theory
\[
    V\hookrightarrow
    \mathbf{Beh}^{\mathrm{can}}_{\mathbb A,\mathbb B}
\]
inherits the restricted canonical Mealy structure and hence an action
\[
    \chi_V:
    \mathbb A^{\mathrm{op}}
    \longrightarrow
    \operatorname{KlEnd}_{\mathbb B}(V).
\]
Its transition--output image is the identity-on-objects regular image
\[
    \operatorname{Op}_{\mathbb A,\mathbb B}(V)
    :=
    \operatorname{RegIm}(\chi_V).
\]
Equivalently, it is the quotient of \(\mathbb A^{\mathrm{op}}\) obtained by
identifying arrows that induce the same transition--output operation on all
states of \(V\).

Fourth, for a behaviour specification
\[
    k:K\to\mathbf{Beh}_0,
\]
the syntactic transition--output category is
\[
    \operatorname{SynTO}_{\mathbb A,\mathbb B}(K)
    :=
    \operatorname{Op}_{\mathbb A,\mathbb B}(\operatorname{Der}_0(K)),
\]
where
\[
    \operatorname{Der}_0(K)\hookrightarrow\mathbf{Beh}_0
\]
is the derivative image constructed in
Chapter~\ref{ch:behaviour-categories-myhill-nerode}. Thus
\(\operatorname{SynTO}(K)\) is the transition--output image of the local
coequational theory generated by the specified behaviours.

The operation image is controlled by the operation congruence
\[
    R_V:=\ker(\chi_V),
\]
an identity-on-objects congruence on \(\mathbb A^{\mathrm{op}}\). Conversely,
an operation congruence
\[
    R\rightrightarrows A_1
\]
determines the largest local coequational theory recognized by the quotient
\[
    \mathbb A^{\mathrm{op}}
    \twoheadrightarrow
    \mathbb A^{\mathrm{op}}/R.
\]
These two constructions form an order-reversing Galois connection. The fixed
points give the local algebraic Eilenberg theorem.

The final part of the chapter globalizes the construction. Input functors act
on residual behaviours by reindexing, local theories are transported by
\(F^\bullet\), and operation congruences are transported by the closure
\[
    F^\flat R:=R_{F^\bullet(V_R)}.
\]
The one-object Boolean regular-language specialization recovers the role of
syntactic monoids and syntactic congruences in Eilenberg-type correspondences
for regular languages
\cite{Eilenberg1974,AdamekMiliusMyersUrbat2019TOCL,Salamanca2017Unveiling}.

\subsection{Standing assumptions and notation}
\label{subsec:syntactic-to-standing-assumptions}

Throughout this chapter we use the assumptions and notation of
Chapter~\ref{ch:coequational-subobjects-birkhoff}. Thus \(\mathcal E\) is a
Grothendieck topos with chosen pullbacks. In particular, \(\mathcal E\) is
Barr-exact and locally cartesian closed, has stable regular epi--mono
factorizations, effective quotients of internal equivalence relations, small
coproducts, and complete subobject lattices in every slice. We use these
properties in the standard topos-theoretic form recorded in
\cite{Johnstone2002,Borceux1994,Jacobs1999CategoricalLogic}.

The input and output internal categories are
\[
    \mathbb A=(A_0,A_1,s_A,t_A,e_A,m_A),
    \qquad
    \mathbb B=(B_0,B_1,s_B,t_B,e_B,m_B).
\]
When \(\mathbb A\) and \(\mathbb B\) are fixed, we write
\[
    Z:=A_0\times B_0.
\]
A Mealy carrier is an object of the slice
\[
    \mathcal E/Z,
\]
displayed as a span
\[
    A_0\xleftarrow{p}P\xrightarrow{q}B_0.
\]

Composition in \(\mathbb A\) is written
\[
    m_A(a,b)=b\circ a
\]
for composable arrows
\[
    a:u\to v,
    \qquad
    b:v\to w.
\]
Thus a state over \(w\) processes \(b\circ a\) by first processing \(b\), then
processing \(a\):
\[
    \delta(b\circ a,x)
    =
    \delta(a,\delta(b,x)).
\]
The corresponding output law is
\[
    \omega(b\circ a,x)
    =
    \omega(b,x)\circ\omega(a,\delta(b,x)).
\]

We regard the arrow object \(A_1\) of
\[
    \mathbb A^{\mathrm{op}}
\]
as an object over
\[
    A_0\times A_0
\]
by
\[
    (t_A,s_A):A_1\to A_0\times A_0.
\]
Thus an arrow
\[
    a:u\to v
\]
of \(\mathbb A\) is read as an arrow
\[
    a^{\mathrm{op}}:v\to u
\]
of \(\mathbb A^{\mathrm{op}}\).

\section{The output-comparison Kleisli category}
\label{sec:to-output-comparison-kleisli-endomorphisms}

\subsection{The output-comparison Kleisli monad}
\label{subsec:to-output-comparison-kleisli-monad}

The output category \(\mathbb B\) induces a monad
\[
    T_{\mathbb B}:\mathcal E/B_0\to\mathcal E/B_0
\]
by
\[
    T_{\mathbb B}:=(t_B)_!(s_B)^*.
\]
This is the fibrewise slice form of the output-comparison monad
\[
    \mathbb B\circ -
\]
used earlier on spans into \(B_0\).

Explicitly, if
\[
    q_X:X\to B_0
\]
is an object of \(\mathcal E/B_0\), then
\[
    T_{\mathbb B}X
    =
    X\times_{q_X,B_0,s_B}B_1
    \xrightarrow{t_B\pi_{B_1}}
    B_0.
\]
A generalized element of \(T_{\mathbb B}X\) is a pair
\[
    (x,\alpha)
\]
where
\[
    \alpha:q_X(x)\to b
\]
is an arrow of \(\mathbb B\). The unit sends
\[
    x\longmapsto (x,e_B(q_Xx)),
\]
and multiplication composes output-comparison arrows in \(\mathbb B\).

The defining pullback is
\[
\begin{tikzcd}[column sep=large, row sep=large]
    T_{\mathbb B}X
        \arrow[r]
        \arrow[d]
        \arrow[dr, phantom, "\lrcorner", very near start]
    &
    B_1
        \arrow[d, "s_B"]
    \\
    X
        \arrow[r, "q_X"']
    &
    B_0 .
\end{tikzcd}
\]
The structure map to \(B_0\) is read after the output comparison:
\[
\begin{tikzcd}[column sep=large, row sep=large]
    T_{\mathbb B}X
        \arrow[r]
        \arrow[dr, "{t_B\pi_{B_1}}"']
    &
    B_1
        \arrow[d, "t_B"]
    \\
    &
    B_0 .
\end{tikzcd}
\]

\begin{proposition}[Kleisli normal form]
\label{prop:to-kleisli-normal-form}
    Let
    \[
        q_X:X\to B_0,
        \qquad
        q_Y:Y\to B_0
    \]
    be objects of \(\mathcal E/B_0\). A Kleisli map
    \[
        X\rightsquigarrow Y
    \]
    for \(T_{\mathbb B}\) is equivalently a pair
    \[
        F:X\to Y,
        \qquad
        \alpha:X\to B_1
    \]
    satisfying
    \[
        s_B\alpha=q_YF,
        \qquad
        t_B\alpha=q_X.
    \]
    Thus
    \[
        \alpha_x:q_Y(Fx)\to q_X(x).
    \]

    If
    \[
        (F,\alpha):X\rightsquigarrow Y,
        \qquad
        (G,\gamma):Y\rightsquigarrow Z
    \]
    are Kleisli maps, their composite is
    \[
        (G,\gamma)\odot(F,\alpha)
        =
        (GF,\alpha\circ_F\gamma),
    \]
    where
    \[
        (\alpha\circ_F\gamma)_x
        =
        \alpha_x\circ\gamma_{Fx}.
    \]
\end{proposition}

\begin{proof}
    \leavevmode
    \begin{enumerate}
        \item \textbf{Unpack a Kleisli map.}
        A Kleisli map
        \[
            X\rightsquigarrow Y
        \]
        is a map
        \[
            X\to T_{\mathbb B}Y
            =
            Y\times_{q_Y,B_0,s_B}B_1
        \]
        in \(\mathcal E/B_0\). Hence it is equivalently a pair
        \[
            F:X\to Y,
            \qquad
            \alpha:X\to B_1
        \]
        satisfying
        \[
            s_B\alpha=q_YF.
        \]

        \item \textbf{Use the slice condition.}
        The condition that
        \[
            X\to T_{\mathbb B}Y
        \]
        be a map over \(B_0\) is
        \[
            t_B\alpha=q_X.
        \]
        Thus \(\alpha\) has type
        \[
            \alpha_x:q_Y(Fx)\to q_X(x).
        \]

        \item \textbf{Compute composition.}
        Applying \((F,\alpha)\) first and \((G,\gamma)\) second gives the
        comparison chain
        \[
        \begin{tikzcd}[column sep=large, row sep=large]
            q_Z(GFx)
                \arrow[r, "\gamma_{Fx}"]
            &
            q_Y(Fx)
                \arrow[r, "\alpha_x"]
            &
            q_X(x).
        \end{tikzcd}
        \]
        Hence the composite comparison arrow is
        \[
            \alpha_x\circ\gamma_{Fx}.
        \]
        This gives the stated composite
        \[
            (G,\gamma)\odot(F,\alpha)
            =
            (GF,\alpha\circ_F\gamma).
        \]
    \end{enumerate}
\end{proof}

\begin{remark}[Direction of output comparison]
\label{rem:to-direction-output-comparison}
    The monad
    \[
        T_{\mathbb B}=(t_B)_!(s_B)^*
    \]
    is chosen so that a Kleisli map
    \[
        X\rightsquigarrow Y
    \]
    carries a comparison arrow
    \[
        q_Y(Fx)\to q_X(x).
    \]
    The comparison arrow points from the new output interface to the old one.
    This matches the Mealy convention
    \[
        \omega(a,x):q(\delta(a,x))\to q(x),
    \]
    where execution of an input arrow moves the state while producing a
    comparison from the new output interface to the old one.
\end{remark}

\subsection{Fibrewise Kleisli endomorphisms}
\label{subsec:to-fibrewise-kleisli-endomorphisms}

Let
\[
    A_0\xleftarrow{p}P\xrightarrow{q}B_0
\]
be an object over
\[
    Z=A_0\times B_0.
\]
For \(v\in A_0\), write \(P_v\) for the fibre of \(p\) over \(v\).

The internal category
\[
    \operatorname{KlEnd}_{\mathbb B}(P)
\]
will have object object \(A_0\). Its arrows \(v\to u\) are Kleisli maps
\[
    P_v\rightsquigarrow P_u.
\]

Put
\[
    A_0^{[2]}:=A_0\times A_0,
\]
and let
\[
    \ell,r:A_0^{[2]}\to A_0
\]
be the first and second projections. A point
\[
    (v,u)\in A_0^{[2]}
\]
is read as the source and target of an arrow
\[
    v\to u.
\]
Define
\[
    P_{\mathrm{src}}
    :=
    A_0^{[2]}\times_{\ell,A_0,p}P,
    \qquad
    P_{\mathrm{tgt}}
    :=
    A_0^{[2]}\times_{r,A_0,p}P.
\]
Over \((v,u)\), these are respectively \(P_v\) and \(P_u\).

The source and target families are displayed by
\[
\begin{tikzcd}[column sep=large, row sep=large]
    P_{\mathrm{src}}
        \arrow[r]
        \arrow[d]
        \arrow[dr, phantom, "\lrcorner", very near start]
    &
    P
        \arrow[d, "p"]
    \\
    A_0^{[2]}
        \arrow[r, "\ell"']
    &
    A_0
\end{tikzcd}
\qquad
\begin{tikzcd}[column sep=large, row sep=large]
    P_{\mathrm{tgt}}
        \arrow[r]
        \arrow[d]
        \arrow[dr, phantom, "\lrcorner", very near start]
    &
    P
        \arrow[d, "p"]
    \\
    A_0^{[2]}
        \arrow[r, "r"']
    &
    A_0 .
\end{tikzcd}
\]

Apply \(T_{\mathbb B}\) fibrewise over \(A_0^{[2]}\):
\[
    T_{\mathbb B}^{[2]}(P_{\mathrm{tgt}})
    :=
    P_{\mathrm{tgt}}\times_{q,B_0,s_B}B_1.
\]
It maps to
\[
    A_0^{[2]}\times B_0
\]
by
\[
    (v,u,y,\alpha)
    \longmapsto
    ((v,u),t_B\alpha).
\]
The object \(P_{\mathrm{src}}\) maps to \(A_0^{[2]}\times B_0\) by
\[
    (v,u,x)
    \longmapsto
    ((v,u),q(x)).
\]
Define
\[
    E_{[2]}
    :=
    P_{\mathrm{src}}
    \times_{A_0^{[2]}\times B_0}
    T_{\mathbb B}^{[2]}(P_{\mathrm{tgt}}).
\]
A point of \(E_{[2]}\) over \((v,u,x)\), where \(x\in P_v\), is a pair
\[
    y\in P_u,
    \qquad
    \alpha:q(y)\to q(x).
\]

The pullback defining \(E_{[2]}\) is
\[
\begin{tikzcd}[column sep=large, row sep=large]
    E_{[2]}
        \arrow[r]
        \arrow[d]
        \arrow[dr, phantom, "\lrcorner", very near start]
    &
    T_{\mathbb B}^{[2]}(P_{\mathrm{tgt}})
        \arrow[d]
    \\
    P_{\mathrm{src}}
        \arrow[r]
    &
    A_0^{[2]}\times B_0 .
\end{tikzcd}
\]

The pullback comes with a projection
\[
    E_{[2]}\to P_{\mathrm{src}}.
\]
It is this object of the slice
\[
    \mathcal E/P_{\mathrm{src}}
\]
to which the dependent product below is applied.

Since \(\mathcal E\) is locally cartesian closed, the dependent product along
\[
    \pi_{\mathrm{src}}:P_{\mathrm{src}}\to A_0^{[2]}
\]
exists. Define
\[
    \operatorname{KlEnd}_{\mathbb B}(P)_1
    :=
    \Pi_{\pi_{\mathrm{src}}}(E_{[2]})
\]
as an object over \(A_0^{[2]}\). It represents fibrewise Kleisli maps
\[
    P_v\rightsquigarrow P_u.
\]
The dependent product comes with evaluation
\[
    \mathrm{ev}:
    \operatorname{KlEnd}_{\mathbb B}(P)_1
    \times_{A_0^{[2]}}
    P_{\mathrm{src}}
    \longrightarrow
    E_{[2]}.
\]
In generalized elements,
\[
    \mathrm{ev}(\varphi,x)
    =
    \bigl(\varphi_0(x),\alpha^\varphi_x\bigr),
\]
where
\[
    \varphi_0(x)\in P_u,
    \qquad
    \alpha^\varphi_x:q(\varphi_0x)\to q(x).
\]

\subsection{The internal category of Kleisli endomorphisms}
\label{subsec:to-klend-internal-category}

\begin{proposition}[Internal category of Kleisli endomorphisms]
\label{prop:to-klend-internal-category}
    The object object \(A_0\) and arrow object
    \[
        \operatorname{KlEnd}_{\mathbb B}(P)_1
    \]
    form an internal category
    \[
        \operatorname{KlEnd}_{\mathbb B}(P).
    \]
    Identities are induced by the unit of \(T_{\mathbb B}\), and composition is
    induced by Kleisli composition for \(T_{\mathbb B}\).
\end{proposition}

\begin{proof}
    \leavevmode
    \begin{enumerate}
        \item \textbf{Define source and target.}
        The source and target maps are the two projections from
        \[
            A_0^{[2]}=A_0\times A_0
        \]
        along the structure map
        \[
            \operatorname{KlEnd}_{\mathbb B}(P)_1\to A_0^{[2]}.
        \]
        A point over \((v,u)\) is read as an arrow
        \[
            v\to u.
        \]

        \item \textbf{Define identities.}
        The identity at \(v\) is the Kleisli identity
        \[
            P_v\rightsquigarrow P_v,
            \qquad
            x\longmapsto (x,e_B(qx)).
        \]
        Internally, this is obtained by transposing the unit section along the
        diagonal
        \[
            A_0\to A_0\times A_0.
        \]

        \item \textbf{Define composable pairs.}
        Let
        \[
            s,t:\operatorname{KlEnd}_{\mathbb B}(P)_1\to A_0
        \]
        be source and target. The object of composable pairs is
        \[
            \operatorname{KlEnd}_{\mathbb B}(P)_2
            :=
            \operatorname{KlEnd}_{\mathbb B}(P)_1
            \times_{t,A_0,s}
            \operatorname{KlEnd}_{\mathbb B}(P)_1.
        \]
        Over a triple \((v,u,w)\), a point is a pair
        \[
            \varphi:P_v\rightsquigarrow P_u,
            \qquad
            \psi:P_u\rightsquigarrow P_w.
        \]

        \item \textbf{Transpose Kleisli composition.}
        Pulling back the evaluation maps along
        \[
            \operatorname{KlEnd}_{\mathbb B}(P)_2
        \]
        gives, for each \(x\in P_v\), the data
        \[
            \varphi(x)=
            \bigl(\varphi_0x,\alpha^\varphi_x\bigr),
            \qquad
            \psi(\varphi_0x)
            =
            \bigl(\psi_0\varphi_0x,\alpha^\psi_{\varphi_0x}\bigr).
        \]
        Kleisli composition gives
        \[
            (\psi\odot\varphi)(x)
            =
            \bigl(
                \psi_0(\varphi_0x),
                \alpha^\varphi_x\circ\alpha^\psi_{\varphi_0x}
            \bigr).
        \]
        The comparison chain is
        \[
        \begin{tikzcd}[column sep=large, row sep=large]
            q(\psi_0(\varphi_0x))
                \arrow[r, "{\alpha^\psi_{\varphi_0x}}"]
            &
            q(\varphi_0x)
                \arrow[r, "{\alpha^\varphi_x}"]
            &
            q(x).
        \end{tikzcd}
        \]
        Transposition along the dependent-product adjunction gives the
        composition map
        \[
            \operatorname{KlEnd}_{\mathbb B}(P)_2
            \to
            \operatorname{KlEnd}_{\mathbb B}(P)_1.
        \]

        \item \textbf{Check the composition convention.}
        With the convention
        \[
            m(c,d)=d\circ c
        \]
        for internal-category composition, the composition just constructed
        satisfies
        \[
            m_{\operatorname{KlEnd}}(\varphi,\psi)
            =
            \psi\odot\varphi.
        \]

        \item \textbf{Verify the internal category laws.}
        Associativity and unitality follow from the monad laws for
        \(T_{\mathbb B}\), equivalently from associativity and unitality in
        \(\mathbb B\). Since all structure maps are obtained by transposition
        from the corresponding fibrewise Kleisli operations, the laws hold
        internally.
    \end{enumerate}
\end{proof}

\begin{lemma}[Regular-epimorphic detection of evaluated Kleisli equality]
\label{lem:to-regular-epi-detects-kleisli-equality}
    Let
    \[
        U\to A_0\times A_0
    \]
    be an object of \(\mathcal E/(A_0\times A_0)\), and let
    \[
        \varphi,\psi:
        U\to \operatorname{KlEnd}_{\mathbb B}(P)_1
    \]
    be maps over \(A_0\times A_0\). Define
    \[
        P_{\mathrm{src},U}
        :=
        U\times_{A_0\times A_0}
        P_{\mathrm{src}},
    \]
    where \(P_{\mathrm{src}}\to A_0\times A_0\) is the source-state family.
    Let
    \[
        e:E\twoheadrightarrow P_{\mathrm{src},U}
    \]
    be a regular epimorphism. If the two evaluated maps
    \[
        E
        \rightrightarrows
        E_{[2]}
    \]
    induced by \(\varphi\) and \(\psi\) agree after precomposition with \(e\),
    then
    \[
        \varphi=\psi.
    \]

    Equivalently, two fibrewise Kleisli operations are equal whenever their
    state components and output-comparison components agree after a
    regular-epimorphic cover of the source-state family.
\end{lemma}

\begin{proof}
    Let
    \[
        \pi_{\mathrm{src},U}:P_{\mathrm{src},U}\to U
    \]
    be the pullback of
    \[
        \pi_{\mathrm{src}}:P_{\mathrm{src}}\to A_0\times A_0.
    \]
    By construction of
    \[
        \operatorname{KlEnd}_{\mathbb B}(P)_1
        =
        \Pi_{\pi_{\mathrm{src}}}(E_{[2]}),
    \]
    and stability of dependent products under pullback, pulling back the representing object along \(U\to A_0\times A_0\) gives the dependent product over \(U\). Hence the maps
    \[
        \varphi,\psi:
        U\to \operatorname{KlEnd}_{\mathbb B}(P)_1
    \]
    correspond, by adjunction, to evaluated maps
    \[
        \widetilde{\varphi},\widetilde{\psi}:
        P_{\mathrm{src},U}\to E_{[2]}
    \]
    over \(P_{\mathrm{src},U}\). Concretely, these are obtained by pulling back
    the universal evaluation map
    \[
        \mathrm{ev}:
        \operatorname{KlEnd}_{\mathbb B}(P)_1
        \times_{A_0\times A_0}
        P_{\mathrm{src}}
        \to
        E_{[2]}.
    \]

    By assumption,
    \[
        \widetilde{\varphi}\,e
        =
        \widetilde{\psi}\,e.
    \]
    Since
    \[
        e:E\twoheadrightarrow P_{\mathrm{src},U}
    \]
    is a regular epimorphism, hence an epimorphism, it follows that
    \[
        \widetilde{\varphi}
        =
        \widetilde{\psi}.
    \]
    Transposing back along the dependent-product adjunction yields
    \[
        \varphi=\psi.
    \]

    The final formulation is simply the same statement unpacked into the state
    and output-comparison components of the evaluated maps into \(E_{[2]}\).
\end{proof}

\section{Mealy morphisms as transition--output actions}
\label{sec:residual-transition-output-actions}

\subsection{Mealy structures as fibrewise Kleisli actions}
\label{subsec:to-mealy-structures-as-kleisli-actions}

\begin{theorem}[Mealy structures as fibrewise Kleisli actions]
\label{thm:to-mealy-structures-as-kleisli-actions}
    Let
    \[
        A_0\xleftarrow{p}P\xrightarrow{q}B_0
    \]
    be a carrier span. To give a Mealy structure on \(P\) is equivalent to
    giving an identity-on-objects internal functor
    \[
        \chi_P:
        \mathbb A^{\mathrm{op}}
        \to
        \operatorname{KlEnd}_{\mathbb B}(P).
    \]
    Under this correspondence, an arrow
    \[
        a:u\to v
    \]
    of \(\mathbb A\), regarded as
    \[
        a^{\mathrm{op}}:v\to u
    \]
    in \(\mathbb A^{\mathrm{op}}\), is sent to the Kleisli map
    \[
        P_v\rightsquigarrow P_u,
        \qquad
        x\longmapsto
        \bigl(\delta(a,x),\omega(a,x)\bigr).
    \]
\end{theorem}

\begin{proof}
    \leavevmode
    \begin{enumerate}
        \item \textbf{Construct the action from a Mealy structure.}
        Given a Mealy structure \((\delta,\omega)\), the typing equations say
        that, for every arrow
        \[
            a:u\to v,
        \]
        the assignment
        \[
            x\longmapsto
            \bigl(\delta(a,x),\omega(a,x)\bigr)
        \]
        is a Kleisli map
        \[
            P_v\rightsquigarrow P_u.
        \]
        By the dependent-product representation of
        \[
            \operatorname{KlEnd}_{\mathbb B}(P)_1,
        \]
        this transposes to an arrow map
        \[
            A_1\to\operatorname{KlEnd}_{\mathbb B}(P)_1
        \]
        over \(A_0\times A_0\), where \(A_1\), as the arrow object of
        \(\mathbb A^{\mathrm{op}}\), lies over \(A_0\times A_0\) by
        \[
            (t_A,s_A):A_1\to A_0\times A_0.
        \]

        \item \textbf{Check identities.}
        The Mealy unit laws say that
        \[
            \delta(e_Av,x)=x,
            \qquad
            \omega(e_Av,x)=e_B(qx).
        \]
        Hence \(\chi_P\) preserves identities.

        \item \textbf{Check composition.}
        Let
        \[
            a:u\to v,
            \qquad
            b:v\to w
        \]
        be arrows of \(\mathbb A\). Then in \(\mathbb A^{\mathrm{op}}\),
        \[
            (b\circ a)^{\mathrm{op}}
            =
            a^{\mathrm{op}}\circ b^{\mathrm{op}}.
        \]
        Therefore functoriality of \(\chi_P\) says
        \[
            \chi_P((b\circ a)^{\mathrm{op}})
            =
            \chi_P(a^{\mathrm{op}})\odot\chi_P(b^{\mathrm{op}}).
        \]
        Suppressing the superscript \({\mathrm{op}}\), this is
        \[
            \chi_P(b\circ a)
            =
            \chi_P(a)\odot\chi_P(b).
        \]

        \item \textbf{Recover the Mealy multiplication laws.}
        Evaluating this Kleisli composite at \(x\in P_w\) gives
        \[
            \delta(b\circ a,x)
            =
            \delta(a,\delta(b,x)),
        \]
        and the comparison chain
        \[
        \begin{tikzcd}[column sep=large, row sep=large]
            q(\delta(a,\delta(b,x)))
                \arrow[r, "{\omega(a,\delta(b,x))}"]
            &
            q(\delta(b,x))
                \arrow[r, "{\omega(b,x)}"]
            &
            q(x).
        \end{tikzcd}
        \]
        Hence
        \[
            \omega(b\circ a,x)
            =
            \omega(b,x)\circ\omega(a,\delta(b,x)).
        \]
        These are precisely the Mealy multiplication laws.

        \item \textbf{Construct the Mealy structure from an action.}
        Conversely, an identity-on-objects internal functor
        \[
            \chi_P:\mathbb A^{\mathrm{op}}
            \to
            \operatorname{KlEnd}_{\mathbb B}(P)
        \]
        evaluates each arrow \(a:u\to v\) and each state \(x\in P_v\) to a
        Kleisli normal form
        \[
            \bigl(\delta(a,x),\omega(a,x)\bigr).
        \]
        The source-target equations of \(\operatorname{KlEnd}\) give the
        Mealy typing equations. Preservation of identities and composition
        gives the Mealy unit and multiplication laws.

        \item \textbf{Verify inverse constructions.}
        The two constructions are inverse by the universal property of the
        dependent product defining
        \[
            \operatorname{KlEnd}_{\mathbb B}(P)_1.
        \]
    \end{enumerate}
\end{proof}

\subsection{The residual action of a local theory}
\label{subsec:to-residual-action-local-theory}

\begin{definition}[Residual transition--output action]
\label{def:to-residual-kleisli-action}
    Let
    \[
        V\hookrightarrow
        \mathbf{Beh}^{\mathrm{can}}_{\mathbb A,\mathbb B}
    \]
    be a local coequational theory. Its \textbf{residual transition--output
    action} is the internal functor
    \[
        \chi_V:
        \mathbb A^{\mathrm{op}}
        \to
        \operatorname{KlEnd}_{\mathbb B}(V)
    \]
    corresponding to the restricted canonical Mealy structure on \(V\).
\end{definition}

Explicitly, for
\[
    a:u\to v,
    \qquad
    H\in V_v,
\]
one has
\[
    \chi_V(a)(H)
    =
    \bigl(\partial_aH,\omega_V(a,H)\bigr),
\]
where \(\chi_V(a)\) denotes the image of \(a^{\mathrm{op}}\) under
\(\chi_V\). The output comparison is the arrow assigned by \(H\) to the
residual morphism
\[
    a\to 1_v
\]
in the slice \(\mathbb A/v\):
\[
\begin{tikzcd}[column sep=large, row sep=large]
    u
        \arrow[r, "a"]
        \arrow[rr, "a", bend left=15]
    &
    v
        \arrow[r, "1_v"']
    &
    v
\end{tikzcd}
\qquad
\leadsto
\qquad
\begin{tikzcd}[column sep=large, row sep=large]
    H(a)
        \arrow[r, "{H(a\to 1_v)}"]
    &
    H(1_v).
\end{tikzcd}
\]

\begin{lemma}[Naturality under strict subobjects]
\label{lem:to-actions-restrict-to-subobjects}
    If
    \[
        V\hookrightarrow W\hookrightarrow
        \mathbf{Beh}^{\mathrm{can}}_{\mathbb A,\mathbb B}
    \]
    are local coequational theories, then the residual action of \(V\) is the
    restriction of the residual action of \(W\).
\end{lemma}

\begin{proof}
    Both \(V\) and \(W\) carry the restricted canonical transition and output
    structure. Since the inclusion
    \[
        V\hookrightarrow W
    \]
    is a strict semantic Mealy homomorphism, the transition square
    \[
    \begin{tikzcd}[column sep=huge, row sep=large]
        A_1\times_{A_0}V
            \arrow[r]
            \arrow[d, "\delta_V"']
        &
        A_1\times_{A_0}W
            \arrow[d, "\delta_W"]
        \\
        V
            \arrow[r, hook]
        &
        W
    \end{tikzcd}
    \]
    commutes, and the output map of \(V\) is the restriction of the output map
    of \(W\):
    \[
    \begin{tikzcd}[column sep=huge, row sep=large]
        A_1\times_{A_0}V
            \arrow[r]
            \arrow[dr, "\omega_V"']
        &
        A_1\times_{A_0}W
            \arrow[d, "\omega_W"]
        \\
        &
        B_1 .
    \end{tikzcd}
    \]
    Hence the induced Kleisli operations agree after restriction to each fibre
    \(V_v\).
\end{proof}

\section{Transition--output images}
\label{sec:transition-output-images}

\subsection{Identity-on-objects regular images}
\label{subsec:to-identity-on-objects-regular-images}

\begin{definition}[Identity-on-objects regular epimorphism]
\label{def:to-regular-epi-identity-on-objects-categories}
    Let
    \[
        F:\mathbb C\to\mathbb D
    \]
    be an internal functor whose object map is the identity on a common object
    object
    \[
        C_0=D_0.
    \]
    We call \(F\) an \textbf{identity-on-objects regular epimorphism} if its
    arrow map is a regular epimorphism in
    \[
        \mathcal E/(C_0\times C_0),
    \]
    where both arrow objects are regarded over \(C_0\times C_0\) by their
    source-target maps. Identity-on-objects monomorphisms are defined
    similarly, replacing regular epimorphism by monomorphism.
\end{definition}

\begin{lemma}[Identity-on-objects regular images]
\label{lem:to-identity-on-objects-regular-images}
    Let
    \[
        F:\mathbb C\to\mathbb D
    \]
    be an identity-on-objects internal functor. Factor its arrow map in
    \[
        \mathcal E/(C_0\times C_0)
    \]
    as
    \[
        C_1\xrightarrow{r}I_1\xrightarrow{m}D_1,
    \]
    where \(r\) is a regular epimorphism and \(m\) is a monomorphism. Then
    \(I_1\), with object object \(C_0\), carries a unique internal category
    structure for which
    \[
        \mathbb C\twoheadrightarrow I\hookrightarrow \mathbb D
    \]
    are internal functors.
\end{lemma}

\begin{proof}
    \leavevmode
    \begin{enumerate}
        \item \textbf{Inherit source and target.}
        The source and target maps of \(I_1\) are inherited from its structure
        as an object over
        \[
            C_0\times C_0.
        \]

        \item \textbf{Define identities.}
        Define the identity map of \(I\) by
        \[
            e_I:=r\,e_{\mathbb C}:C_0\to I_1.
        \]
        Since \(F\) preserves identities and \(F_1=mr\), one has
        \[
            me_I
            =
            mr e_{\mathbb C}
            =
            F_1e_{\mathbb C}
            =
            e_{\mathbb D}.
        \]
        Thus the identity map of \(I\) is the unique map whose composite with
        \(m\) is the identity map of \(\mathbb D\).

        \item \textbf{Construct the cover of composable pairs.}
        Let
        \[
            C_2=C_1\times_{C_0}C_1,
            \qquad
            I_2=I_1\times_{C_0}I_1
        \]
        be the objects of composable pairs. The induced map
        \[
            C_2\to I_2
        \]
        is a regular epimorphism, by stability of regular epimorphisms under
        pullback and finite products in the relevant topos slices.

        \item \textbf{Descend composition.}
        Define
        \[
            h:I_2\to D_1
        \]
        by composing the two image arrows in \(\mathbb D\). This is a map
        over \(C_0\times C_0\), where \(I_2\) is regarded over
        \(C_0\times C_0\) by the source of the first arrow and the target
        of the second arrow. After
        precomposition with
        \[
            C_2\twoheadrightarrow I_2,
        \]
        this map is
        \[
            m_D(F_1c,F_1d)
            =
            F_1(m_C(c,d))
            =
            mr(m_C(c,d)).
        \]
        Hence \(h\) lands in the subobject
        \[
            m:I_1\hookrightarrow D_1
        \]
        after a regular-epimorphic cover. Regular-epimorphic descent for
        subobject membership gives a unique factorization
        \[
            I_2\to I_1.
        \]
        This factorization is the composition map of \(I\).

        \item \textbf{Verify associativity and units.}
        For associativity, use the induced regular epimorphism
        \[
            C_3\twoheadrightarrow I_3
        \]
        on composable triples. After precomposition with this cover, the two
        associativity composites in \(I\) become the corresponding composites
        in \(\mathbb D\), hence agree. Since the cover is epic, associativity
        holds in \(I\).

        The unit laws are checked in the same way using
        \[
            C_1\twoheadrightarrow I_1.
        \]

        \item \textbf{Prove uniqueness.}
        Any internal category structure on \(I_1\) making
        \[
            m:I\hookrightarrow\mathbb D
        \]
        a functor must have the identity and composition just constructed.
        Since \(m\) is monic, such a structure is unique.
    \end{enumerate}
\end{proof}

\begin{lemma}[Identity-on-objects quotient by a congruence]
\label{lem:to-identity-on-objects-congruence-quotient}
    Let
    \[
        R\rightrightarrows C_1
    \]
    be an identity-on-objects operation congruence on an internal category
    \[
        \mathbb C.
    \]
    Then the effective quotient
    \[
        C_1\twoheadrightarrow C_1/R
    \]
    in
    \[
        \mathcal E/(C_0\times C_0)
    \]
    carries a unique internal category structure with object object \(C_0\)
    such that
    \[
        \mathbb C\twoheadrightarrow \mathbb C/R
    \]
    is an identity-on-objects internal functor.
\end{lemma}

\begin{proof}
    Identities descend because \(R\) is compatible with identities.
    Composition descends because \(R\) is compatible with composition. The
    internal category laws hold after precomposition with the regular
    epimorphic quotient cover, where they become the laws of \(\mathbb C\).
    Uniqueness follows because the quotient map is epic on arrows.
\end{proof}

\subsection{Transition--output images}
\label{subsec:to-operation-images}

\begin{definition}[Transition--output image]
\label{def:to-operation-image}
    Let
    \[
        V\hookrightarrow
        \mathbf{Beh}^{\mathrm{can}}_{\mathbb A,\mathbb B}
    \]
    be a local coequational theory. Its \textbf{transition--output image} is
    the identity-on-objects regular image of
    \[
        \chi_V:
        \mathbb A^{\mathrm{op}}
        \to
        \operatorname{KlEnd}_{\mathbb B}(V).
    \]
    It is denoted
    \[
        \operatorname{Op}_{\mathbb A,\mathbb B}(V),
    \]
    or simply \(\operatorname{Op}(V)\).
\end{definition}

Thus \(\operatorname{Op}(V)\) has object object \(A_0\), and its arrow object
is the regular image of
\[
    A_1\to \operatorname{KlEnd}_{\mathbb B}(V)_1
\]
in the slice
\[
    \mathcal E/(A_0\times A_0),
\]
where \(A_1\), as the arrow object of \(\mathbb A^{\mathrm{op}}\), lies over
\(A_0\times A_0\) by
\[
    (t_A,s_A).
\]
The defining image factorization is
\[
\begin{tikzcd}[column sep=large, row sep=large]
    A_1
        \arrow[rr, "\chi_{V,1}"]
        \arrow[dr, two heads]
    &&
    \operatorname{KlEnd}_{\mathbb B}(V)_1
    \\
    &
    \operatorname{Op}(V)_1
        \arrow[ur, hook]
    &
\end{tikzcd}
\qquad
\text{in }\mathcal E/(A_0\times A_0).
\]

\begin{proposition}[Image subcategory]
\label{prop:to-operation-image-subcategory}
    The transition--output image
    \[
        \operatorname{Op}(V)
    \]
    is an internal subcategory of
    \[
        \operatorname{KlEnd}_{\mathbb B}(V).
    \]
\end{proposition}

\begin{proof}
    The residual action
    \[
        \chi_V:
        \mathbb A^{\mathrm{op}}
        \to
        \operatorname{KlEnd}_{\mathbb B}(V)
    \]
    is identity-on-objects. Lemma~\ref{lem:to-identity-on-objects-regular-images}
    creates the regular image as an internal category with object object
    \(A_0\). The monic part of the image factorization is an identity-on-
    objects internal functor
    \[
        \operatorname{Op}(V)\hookrightarrow
        \operatorname{KlEnd}_{\mathbb B}(V),
    \]
    so \(\operatorname{Op}(V)\) is an internal subcategory.
\end{proof}

\subsection{Operation congruences}
\label{subsec:to-operation-congruences}

\begin{definition}[Operation congruence of a local theory]
\label{def:to-operation-congruence-of-local-theory}
    Let
    \[
        V\hookrightarrow
        \mathbf{Beh}^{\mathrm{can}}_{\mathbb A,\mathbb B}
    \]
    be a local coequational theory. Its \textbf{operation congruence}
    \[
        R_V
    \]
    is the kernel pair of
    \[
        \chi_{V,1}:A_1\to
        \operatorname{KlEnd}_{\mathbb B}(V)_1
    \]
    in
    \[
        \mathcal E/(A_0\times A_0),
    \]
    with \(A_1\) regarded over \(A_0\times A_0\) by
    \[
        (t_A,s_A).
    \]
\end{definition}

Thus, for parallel arrows
\[
    a,a':u\to v
\]
of \(\mathbb A\), equivalently arrows lying over the same point
\[
    (v,u)\in A_0\times A_0
\]
as arrows of \(\mathbb A^{\mathrm{op}}\), the relation
\[
    a\,R_V\,a'
\]
means that \(a\) and \(a'\) induce the same Kleisli map
\[
    V_v\rightsquigarrow V_u.
\]
Equivalently, for every \(H\in V_v\),
\[
    \partial_aH=\partial_{a'}H
\]
and
\[
    \omega_V(a,H)=\omega_V(a',H),
\]
where the second equality is typed using the first equality.

\begin{lemma}[Kernel operation congruence]
\label{lem:to-kernel-operation-congruence}
    For every local coequational theory
    \[
        V\hookrightarrow
        \mathbf{Beh}^{\mathrm{can}}_{\mathbb A,\mathbb B},
    \]
    the relation
    \[
        R_V=\ker(\chi_{V,1})
    \]
    is an identity-on-objects operation congruence on
    \(\mathbb A^{\mathrm{op}}\).
\end{lemma}

\begin{proof}
    Since \(R_V\) is the kernel pair of a morphism in the exact slice
    \[
        \mathcal E/(A_0\times A_0),
    \]
    it is an effective internal equivalence relation over
    \(A_0\times A_0\).

    Since
    \[
        \chi_V:
        \mathbb A^{\mathrm{op}}\to\operatorname{KlEnd}_{\mathbb B}(V)
    \]
    is an internal functor, its kernel pair in the arrow slice is stable under
    identity and composition. Thus \(R_V\rightrightarrows A_1\) is not merely
    an internal equivalence relation over \(A_0\times A_0\), but an internal
    congruence on \(\mathbb A^{\mathrm{op}}\).
\end{proof}

\begin{theorem}[Transition--output image quotient]
\label{thm:to-operation-image-quotient}
    There is a canonical isomorphism of identity-on-objects internal categories
    \[
        \operatorname{Op}(V)
        \cong
        \mathbb A^{\mathrm{op}}/R_V.
    \]
\end{theorem}

\begin{proof}
    By Lemma~\ref{lem:to-identity-on-objects-congruence-quotient},
    \[
        \mathbb A^{\mathrm{op}}/R_V
    \]
    exists. Since
    \[
        R_V=\ker(\chi_{V,1}),
    \]
    exactness identifies the quotient
    \[
        A_1/R_V
    \]
    with the regular image of
    \[
        \chi_{V,1}:A_1\to\operatorname{KlEnd}_{\mathbb B}(V)_1
    \]
    in
    \[
        \mathcal E/(A_0\times A_0).
    \]
    The internal category structure on this regular image is the one created
    by Lemma~\ref{lem:to-identity-on-objects-regular-images}, and the quotient
    category structure is unique by
    Lemma~\ref{lem:to-identity-on-objects-congruence-quotient}. Therefore the
    two identity-on-objects internal categories are canonically isomorphic.
\end{proof}

\begin{proposition}[Contravariance under inclusion]
\label{prop:to-operation-image-contravariant}
    Let
    \[
        V\hookrightarrow W\hookrightarrow
        \mathbf{Beh}^{\mathrm{can}}_{\mathbb A,\mathbb B}
    \]
    be local coequational theories. Then
    \[
        R_W\leq R_V
    \]
    and there is a canonical identity-on-objects regular epimorphism
    \[
        \operatorname{Op}(W)\twoheadrightarrow \operatorname{Op}(V).
    \]
\end{proposition}

\begin{proof}
    If two arrows induce the same transition--output operation on every state
    of \(W\), then they induce the same transition--output operation on every
    state of the subobject \(V\). Hence
    \[
        R_W\leq R_V.
    \]
    Since both congruences are effective and
    \[
        R_W\leq R_V,
    \]
    the universal property of the effective quotient by \(R_W\) gives a unique
    identity-on-objects functor
    \[
        \mathbb A^{\mathrm{op}}/R_W
        \to
        \mathbb A^{\mathrm{op}}/R_V.
    \]
    Its arrow map is a regular epimorphism, since it is the comparison between
    effective quotients associated to an inclusion of equivalence relations in
    an exact category. Translating through
    Theorem~\ref{thm:to-operation-image-quotient} gives the desired map
    \[
        \operatorname{Op}(W)\twoheadrightarrow \operatorname{Op}(V).
    \]
\end{proof}

\section{Syntactic transition--output categories}
\label{sec:syntactic-transition-output-categories}

\subsection{Definition}
\label{subsec:to-syntactic-categories-definition}

Let
\[
    k:K\to\mathbf{Beh}_0
\]
be a behaviour specification, and let
\[
    \operatorname{Der}_0(K)\hookrightarrow\mathbf{Beh}_0
\]
be its derivative image.

\begin{definition}[Syntactic transition--output category]
\label{def:to-syntactic-transition-output-category}
    The \textbf{syntactic transition--output category} of \(K\) is
    \[
        \operatorname{SynTO}_{\mathbb A,\mathbb B}(K)
        :=
        \operatorname{Op}_{\mathbb A,\mathbb B}(\operatorname{Der}_0(K)).
    \]
    Equivalently,
    \[
        \operatorname{SynTO}_{\mathbb A,\mathbb B}(K)
        =
        \operatorname{RegIm}
        \left(
            \mathbb A^{\mathrm{op}}
            \to
            \operatorname{KlEnd}_{\mathbb B}(\operatorname{Der}_0(K))
        \right),
    \]
    with regular image taken on arrow objects over \(A_0\times A_0\).
\end{definition}

\begin{definition}[Syntactic operation congruence]
\label{def:to-syntactic-operation-congruence}
    For parallel arrows
    \[
        a,a':u\to v
    \]
    of \(\mathbb A\), write
    \[
        a\equiv^{\mathrm{syn}}_K a'
    \]
    when
    \[
        a\,R_{\operatorname{Der}_0(K)}\,a'.
    \]
\end{definition}

Thus
\[
    \operatorname{SynTO}_{\mathbb A,\mathbb B}(K)
    \cong
    \mathbb A^{\mathrm{op}}/
    R_{\operatorname{Der}_0(K)}.
\]
The defining quotient is displayed by
\[
\begin{tikzcd}[column sep=large, row sep=large]
    \mathbb A^{\mathrm{op}}
        \arrow[rr]
        \arrow[dr, two heads]
    &&
    \operatorname{KlEnd}_{\mathbb B}(\operatorname{Der}_0(K))
    \\
    &
    \operatorname{SynTO}_{\mathbb A,\mathbb B}(K)
        \arrow[ur, hook]
    &
\end{tikzcd}
\]
where the diagram is identity-on-objects and the image is taken on arrows.

\subsection{Operation images of realizations}
\label{subsec:to-operation-images-mealy-morphisms}

\begin{definition}[Transition--output image of a Mealy morphism]
\label{def:to-operation-image-mealy-morphism}
    For a Mealy morphism \(P\), define
    \[
        \operatorname{Op}(P)
        :=
        \operatorname{Op}(\operatorname{Beh}(P)).
    \]
\end{definition}

\begin{theorem}[Realization quotient]
\label{thm:to-mealy-operation-quotient}
    Let
    \[
        i:K\to P
    \]
    be a realization of
    \[
        k:K\to\mathbf{Beh}_0,
    \]
    so that
    \[
        \beta_P i=k.
    \]
    Then there is a canonical identity-on-objects regular epimorphism
    \[
        \operatorname{Op}(P)
        \twoheadrightarrow
        \operatorname{SynTO}_{\mathbb A,\mathbb B}(K).
    \]
    If \(P\) is reachable as a realization of \(K\), then this regular
    epimorphism is an isomorphism.
\end{theorem}

\begin{proof}
    Since
    \[
        k=\beta_Pi,
    \]
    the specified behaviours lie in the semantic image
    \[
        \operatorname{Beh}(P)\hookrightarrow
        \mathbf{Beh}^{\mathrm{can}}.
    \]
    Because \(\operatorname{Beh}(P)\) is a local coequational theory, it is
    derivative closed. Since it contains the image of \(k\), the derivative
    image of \(K\) factors through it:
    \[
        \operatorname{Der}_0(K)\hookrightarrow\operatorname{Beh}(P).
    \]
    By Proposition~\ref{prop:to-operation-image-contravariant}, this inclusion
    induces a regular epimorphism
    \[
        \operatorname{Op}(\operatorname{Beh}(P))
        \twoheadrightarrow
        \operatorname{Op}(\operatorname{Der}_0(K)).
    \]
    By definition, this is
    \[
        \operatorname{Op}(P)
        \twoheadrightarrow
        \operatorname{SynTO}_{\mathbb A,\mathbb B}(K).
    \]

    Suppose now that \(P\) is reachable. Then the reachability map
    \[
        \rho_{P,K}:D_K\to P
    \]
    is a regular epimorphism and satisfies
    \[
        \beta_P\rho_{P,K}=d_K.
    \]
    Hence the regular images of
    \[
        d_K:D_K\to\mathbf{Beh}_0
        \qquad\text{and}\qquad
        \beta_P:P\to\mathbf{Beh}_0
    \]
    are canonically isomorphic. Therefore
    \[
        \operatorname{Beh}(P)\cong \operatorname{Der}_0(K)
    \]
    as restricted canonical objects. Applying \(\operatorname{Op}\) gives
    \[
        \operatorname{Op}(P)
        \cong
        \operatorname{SynTO}_{\mathbb A,\mathbb B}(K).
    \]
\end{proof}

\begin{remark}[Separated realizations]
\label{rem:to-separated-realizations-operation-image}
    Behavioural separatedness is not used in
    Theorem~\ref{thm:to-mealy-operation-quotient}. Reachability identifies the
    semantic image of \(P\) with \(\operatorname{Der}_0(K)\). If \(P\) is also
    behaviourally separated, the minimal-realization theorem identifies \(P\)
    itself with \(\mathsf{Min}(K)\).
\end{remark}

\subsection{Context expansion}
\label{subsec:to-context-expansion}

For a generalized element
\[
    \kappa\in K
\]
and a residual arrow
\[
    r:v\to p_K(\kappa),
\]
if
\[
    a:u\to v,
\]
then
\[
    r\circ a:u\to p_K(\kappa)
\]
is an object of the residual slice \(\mathbb A/p_K(\kappa)\), and the
triangle
\[
    r\circ a\to r
\]
is a morphism of that slice. We write
\[
    k(\kappa)(a,r)
    :=
    k(\kappa)(r\circ a\to r)
\]
for the arrow of \(\mathbb B\) assigned by the residual behaviour \(k(\kappa)\)
to the residual morphism
\[
    r\circ a\to r
\]
in \(\mathbb A/p_K(\kappa)\).

\begin{theorem}[Context expansion]
\label{thm:to-context-expansion}
    Let
    \[
        a,a':u\to v
    \]
    be parallel arrows of \(\mathbb A\). Then
    \[
        a\equiv^{\mathrm{syn}}_K a'
    \]
    if and only if the following two conditions hold on the regular-epimorphic
    cover of \(\operatorname{Der}_0(K)_v\) by generated residual states.

    For every generalized element \(\kappa\in K\) and every residual context
    \[
        r:v\to p_K(\kappa),
    \]
    one has
    \[
        \partial_{r\circ a}k(\kappa)
        =
        \partial_{r\circ a'}k(\kappa),
    \]
    and
    \[
        k(\kappa)(a,r)=k(\kappa)(a',r).
    \]
    The second equality is equality in the fibre of \(B_1\) over the common
    pair
    \[
    \bigl(
        q(\partial_{r\circ a}k(\kappa)),
        q(\partial_rk(\kappa))
    \bigr)
    =
    \bigl(
        q(\partial_{r\circ a'}k(\kappa)),
        q(\partial_rk(\kappa))
    \bigr).
    \]
\end{theorem}

\begin{proof}
    The derivative image is the regular image of
    \[
        d_K:D_K\to\mathbf{Beh}_0.
    \]
    Pulling back the regular epimorphism
    \[
        D_K\twoheadrightarrow \operatorname{Der}_0(K)
    \]
    along the fibre inclusion
    \[
        \operatorname{Der}_0(K)_v
        \hookrightarrow
        \operatorname{Der}_0(K)
    \]
    gives a regular-epimorphic cover
    \[
        D_K\times_{\operatorname{Der}_0(K)}
        \operatorname{Der}_0(K)_v
        \twoheadrightarrow
        \operatorname{Der}_0(K)_v.
    \]
    Generalized elements of this cover are pairs
    \[
        (r,\kappa),
        \qquad
        r:v\to p_K(\kappa),
    \]
    representing the generated residual state
    \[
        \partial_r k(\kappa).
    \]

    Acting by
    \[
        a:u\to v
    \]
    on this state gives
    \[
        \partial_a(\partial_r k(\kappa))
        =
        \partial_{r\circ a}k(\kappa).
    \]
    Similarly, acting by \(a'\) gives
    \[
        \partial_{r\circ a'}k(\kappa).
    \]
    The output comparison produced at the state
    \[
        \partial_r k(\kappa)
    \]
    along \(a\) is the arrow assigned by \(k(\kappa)\) to the residual
    morphism
    \[
        r\circ a\to r,
    \]
    namely
    \[
        k(\kappa)(a,r).
    \]
    The analogous output for \(a'\) is
    \[
        k(\kappa)(a',r).
    \]

    Equality of the two Kleisli maps
    \[
        \operatorname{Der}_0(K)_v
        \rightsquigarrow
        \operatorname{Der}_0(K)_u
    \]
    is equality after evaluation on the fibre
    \[
        \operatorname{Der}_0(K)_v.
    \]
    Since the displayed map onto this fibre is a regular epimorphism, equality
    is detected after precomposition with that cover by
    Lemma~\ref{lem:to-regular-epi-detects-kleisli-equality}. Therefore the two
    Kleisli maps are equal precisely when the transition components and
    output-comparison components agree on all generated residual states.
\end{proof}


\subsection{Recognition by syntactic quotients}
\label{subsec:to-recognition-by-syntactic-quotients}

The syntactic transition--output category is characterized by a universal
property among identity-on-objects operation quotients of \(\mathbb A^{\mathrm{op}}\).

\subsection{Recognition of specifications}
\label{subsec:to-recognition-specifications}

\begin{definition}[Recognition of a specification]
\label{def:to-operation-quotient-recognizes-specification}
    Let
    \[
        k:K\to\mathbf{Beh}_0
    \]
    be a behaviour specification. An operation quotient
    \[
        q_R:
        \mathbb A^{\mathrm{op}}
        \twoheadrightarrow
        \mathbb A^{\mathrm{op}}/R
    \]
    \textbf{recognizes} \(K\) if it recognizes the derivative theory
    \[
        \operatorname{Der}_0(K)\hookrightarrow\mathbf{Beh}_0.
    \]
\end{definition}

Thus recognition of \(K\) is recognition of all residual behaviours generated
by \(K\), not merely the initial specified behaviours.

\subsection{The syntactic quotient}
\label{subsec:to-syntactic-universal-property}

\begin{theorem}[Syntactic transition--output universal property]
\label{thm:to-syntactic-universal-property}
    Let
    \[
        k:K\to\mathbf{Beh}_0
    \]
    be a behaviour specification, and let
    \[
        q_R:
        \mathbb A^{\mathrm{op}}
        \twoheadrightarrow
        \mathbb A^{\mathrm{op}}/R
    \]
    be an identity-on-objects operation quotient. The following are equivalent:
    \begin{enumerate}
        \item \(q_R\) recognizes \(K\);
        \item \(q_R\) recognizes \(\operatorname{Der}_0(K)\);
        \item
        \[
            R\leq R_{\operatorname{Der}_0(K)};
        \]
        \item there is a canonical identity-on-objects regular epimorphism
        \[
            \mathbb A^{\mathrm{op}}/R
            \twoheadrightarrow
            \operatorname{SynTO}_{\mathbb A,\mathbb B}(K).
        \]
    \end{enumerate}
\end{theorem}

\begin{proof}
    The equivalence of (1) and (2) is definitional. The equivalence of (2) and
    (3) is Lemma~\ref{lem:to-operation-quotient-recognition-criterion}. By
    Theorem~\ref{thm:to-operation-image-quotient},
    \[
        \operatorname{SynTO}_{\mathbb A,\mathbb B}(K)
        =
        \operatorname{Op}(\operatorname{Der}_0(K))
        \cong
        \mathbb A^{\mathrm{op}}/
        R_{\operatorname{Der}_0(K)}.
    \]
    Since identity-on-objects quotients are ordered by inclusion of kernel
    congruences,
    \[
        R\leq R_{\operatorname{Der}_0(K)}
    \]
    is equivalent to a canonical regular epimorphism
    \[
        \mathbb A^{\mathrm{op}}/R
        \twoheadrightarrow
        \mathbb A^{\mathrm{op}}/
        R_{\operatorname{Der}_0(K)}
        \cong
        \operatorname{SynTO}_{\mathbb A,\mathbb B}(K).
    \]
\end{proof}

\begin{corollary}[Syntactic transition--output quotient]
\label{cor:to-syntactic-transition-output-quotient}
    The syntactic transition--output category
    \[
        \operatorname{SynTO}_{\mathbb A,\mathbb B}(K)
    \]
    is the coarsest transition--output quotient of
    \[
        \mathbb A^{\mathrm{op}}
    \]
    recognizing \(K\). Equivalently, every recognizing quotient maps
    canonically and regularly epimorphically onto it.
\end{corollary}

\section{Local algebraic Eilenberg duality}
\label{sec:local-algebraic-eilenberg-duality}

\subsection{Operation quotients and recognition}
\label{subsec:to-operation-quotients-recognition}

The quotient
\[
    \operatorname{Op}(V)
\]
is controlled by the operation congruence \(R_V\). The converse construction
assigns to an operation congruence the largest local coequational theory whose
transition--output action factors through the corresponding quotient.

\begin{definition}[Identity-on-objects operation congruence]
\label{def:to-identity-on-objects-operation-congruence}
    An \textbf{identity-on-objects operation congruence} on
    \(\mathbb A^{\mathrm{op}}\) is an effective internal equivalence relation
    \[
        R\rightrightarrows A_1
    \]
    in the slice
    \[
        \mathcal E/(A_0\times A_0),
    \]
    where \(A_1\) is regarded over \(A_0\times A_0\) by
    \[
        (t_A,s_A),
    \]
    such that \(R\) is compatible with identities and composition.

    Equivalently, \(R\) is the kernel pair of an identity-on-objects regular
    quotient
    \[
        q_R:
        \mathbb A^{\mathrm{op}}
        \twoheadrightarrow
        \mathbb A^{\mathrm{op}}/R.
    \]
    We write
    \[
        \mathsf{Cong}_{\mathrm{ioo}}(\mathbb A^{\mathrm{op}})
    \]
    for the poset of such congruences, ordered by inclusion.
\end{definition}

Thus
\[
    R\leq S
\]
means that every pair identified by \(R\) is identified by \(S\). Equivalently,
the quotient by \(R\) maps regularly epimorphically onto the quotient by \(S\):
\[
\begin{tikzcd}[column sep=large, row sep=large]
    \mathbb A^{\mathrm{op}}
        \arrow[r, two heads]
        \arrow[dr, two heads]
    &
    \mathbb A^{\mathrm{op}}/R
        \arrow[d, two heads, dashed]
    \\
    &
    \mathbb A^{\mathrm{op}}/S .
\end{tikzcd}
\]

\begin{definition}[Recognition by an operation quotient]
\label{def:to-recognition-by-operation-quotient}
    Let
    \[
        q_R:
        \mathbb A^{\mathrm{op}}
        \twoheadrightarrow
        \mathbb A^{\mathrm{op}}/R
    \]
    be an operation quotient. We say that \(q_R\) \textbf{recognizes} a local
    coequational theory
    \[
        V\hookrightarrow\mathbf{Beh}^{\mathrm{can}}
    \]
    if the action
    \[
        \chi_V:
        \mathbb A^{\mathrm{op}}\to
        \operatorname{KlEnd}_{\mathbb B}(V)
    \]
    factors through \(q_R\):
    \[
    \begin{tikzcd}[column sep=large, row sep=large]
        \mathbb A^{\mathrm{op}}
            \arrow[rr, "\chi_V"]
            \arrow[dr, two heads, "q_R"']
        &&
        \operatorname{KlEnd}_{\mathbb B}(V)
        \\
        &
        \mathbb A^{\mathrm{op}}/R
            \arrow[ur, dashed]
        &
    \end{tikzcd}
    \]
\end{definition}

\begin{lemma}[Recognition criterion]
\label{lem:to-operation-quotient-recognition-criterion}
    An operation quotient
    \[
        q_R:
        \mathbb A^{\mathrm{op}}
        \twoheadrightarrow
        \mathbb A^{\mathrm{op}}/R
    \]
    recognizes \(V\) if and only if
    \[
        R\leq R_V.
    \]
\end{lemma}

\begin{proof}
    The action \(\chi_V\) factors through \(q_R\) precisely when every pair of
    arrows identified by \(R\) has the same image under \(\chi_V\). The pairs
    of arrows with the same image under \(\chi_V\) form the kernel pair
    \[
        R_V=\ker(\chi_{V,1}).
    \]
    Hence factorization through \(q_R\) is equivalent to
    \[
        R\leq R_V.
    \]
\end{proof}

\subsection{The largest theory recognized by a congruence}
\label{subsec:to-largest-recognized-theory}

\begin{definition}[Largest theory recognized by a congruence]
\label{def:to-theory-recognized-by-congruence}
    For
    \[
        R\in
        \mathsf{Cong}_{\mathrm{ioo}}(\mathbb A^{\mathrm{op}}),
    \]
    define
    \[
        V_R
        :=
        \bigvee
        \left\{
            S\hookrightarrow\mathbf{Beh}^{\mathrm{can}}
            \mid
            S \text{ is a local coequational theory and } R\leq R_S
        \right\}.
    \]
    The join is taken over a chosen small representative set of the indicated
    subobjects of \(\mathbf{Beh}_0\) in \(\mathcal E/Z\).
\end{definition}

\begin{proposition}[Largest recognized theory]
\label{prop:to-largest-recognized-theory}
    The subobject \(V_R\) is a local coequational theory, satisfies
    \[
        R\leq R_{V_R},
    \]
    and is largest among local coequational theories recognized by \(q_R\).
\end{proposition}

\begin{proof}
    Local coequational theories are closed under small joins. Therefore
    \(V_R\) is a local coequational theory.

    Represent \(V_R\) as the regular image of the coproduct of all local
    theories \(S\) satisfying
    \[
        R\leq R_S:
    \]
    \[
    \begin{tikzcd}[column sep=large, row sep=large]
        \coprod_S S
            \arrow[rr]
            \arrow[dr, two heads]
        &&
        \mathbf{Beh}^{\mathrm{can}}
        \\
        &
        V_R
            \arrow[ur, hook]
        &
    \end{tikzcd}
    \]
    in the strict Mealy category, equivalently on carriers in \(\mathcal E/Z\).

    Equivalently, work internally over the relation object \(R\). Let
    \[
        a,a':u\to v
    \]
    represent a pair classified by \(R\). Pull back the regular epimorphism
    \[
        \coprod_S S\twoheadrightarrow V_R
    \]
    along the fibre
    \[
        (V_R)_v\hookrightarrow V_R.
    \]
    This gives a regular epimorphism
    \[
        E_v\twoheadrightarrow (V_R)_v.
    \]
    On each summand \(S\), the assumption \(R\leq R_S\) says that \(a\) and
    \(a'\) induce the same Kleisli operation
    \[
        S_v\rightsquigarrow S_u.
    \]
    Thus, after precomposition with \(E_v\), the state components agree and the
    output-comparison components agree.

    By Lemma~\ref{lem:to-regular-epi-detects-kleisli-equality}, equality of
    the evaluated Kleisli operations descends along
    \[
        E_v\twoheadrightarrow (V_R)_v.
    \]
    Therefore \(a\) and \(a'\) induce the same Kleisli operation on
    \((V_R)_v\), so
    \[
        R\leq R_{V_R}.
    \]

    Finally, if \(S\) is any local theory recognized by \(q_R\), then
    \[
        R\leq R_S,
    \]
    so \(S\) occurs in the defining join of \(V_R\). Hence
    \[
        S\leq V_R.
    \]
\end{proof}

\begin{lemma}[Order reversal for recognized theories]
\label{lem:to-VR-order-reversing}
    If
    \[
        R\leq S
    \]
    are operation congruences on \(\mathbb A^{\mathrm{op}}\), then
    \[
        V_S\leq V_R.
    \]
\end{lemma}

\begin{proof}
    If a local coequational theory \(T\) is recognized by \(S\), then
    \[
        S\leq R_T.
    \]
    Since
    \[
        R\leq S,
    \]
    one also has
    \[
        R\leq R_T.
    \]
    Hence every summand appearing in the defining join of \(V_S\) also appears
    among the summands defining \(V_R\). Therefore
    \[
        V_S\leq V_R.
    \]
\end{proof}

\subsection{The local operation Galois theorem}
\label{subsec:to-local-operation-galois}

\begin{theorem}[Local operation Galois theorem]
\label{thm:to-local-operation-galois}
    For every local coequational theory \(V\) and every operation congruence
    \(R\), one has
    \[
        V\leq V_R
        \quad\Longleftrightarrow\quad
        R\leq R_V.
    \]
\end{theorem}

\begin{proof}
    Suppose
    \[
        V\leq V_R.
    \]
    By contravariance under inclusion,
    \[
        R_{V_R}\leq R_V.
    \]
    By Proposition~\ref{prop:to-largest-recognized-theory},
    \[
        R\leq R_{V_R}.
    \]
    Hence
    \[
        R\leq R_V.
    \]

    Conversely, suppose
    \[
        R\leq R_V.
    \]
    Then \(V\) appears among the local theories in the defining join of
    \(V_R\). Hence
    \[
        V\leq V_R.
    \]
\end{proof}

\begin{definition}[Saturation and closure]
\label{def:to-operation-saturated-and-closed}
    A local coequational theory \(V\) is \textbf{operation-saturated} if
    \[
        V=V_{R_V}.
    \]
    An operation congruence \(R\) is \textbf{operation-closed} if
    \[
        R=R_{V_R}.
    \]
\end{definition}

\begin{lemma}[The induced closure operators]
\label{lem:to-operation-closure-operators}
    The composites
    \[
        V\longmapsto V_{R_V},
        \qquad
        R\longmapsto R_{V_R}
    \]
    are inflationary and idempotent. More explicitly,
    \[
        V\leq V_{R_V},
        \qquad
        R\leq R_{V_R},
    \]
    and
    \[
        R_{V_{R_V}}=R_V,
        \qquad
        V_{R_{V_R}}=V_R.
    \]
\end{lemma}

\begin{proof}
    The inequality
    \[
        V\leq V_{R_V}
    \]
    follows from Theorem~\ref{thm:to-local-operation-galois} applied to
    \[
        R_V\leq R_V.
    \]
    The inequality
    \[
        R\leq R_{V_R}
    \]
    is Proposition~\ref{prop:to-largest-recognized-theory}.

    Proposition~\ref{prop:to-largest-recognized-theory} gives
    \[
        R_V\leq R_{V_{R_V}}.
    \]
    Since
    \[
        V\leq V_{R_V},
    \]
    contravariance gives
    \[
        R_{V_{R_V}}\leq R_V.
    \]
    Hence
    \[
        R_{V_{R_V}}=R_V.
    \]

    The local Galois theorem gives
    \[
        V_R\leq V_{R_{V_R}}
    \]
    from
    \[
        R_{V_R}\leq R_{V_R}.
    \]
    Conversely, since
    \[
        R\leq R_{V_R},
    \]
    Lemma~\ref{lem:to-VR-order-reversing} gives
    \[
        V_{R_{V_R}}\leq V_R.
    \]
    Therefore
    \[
        V_{R_{V_R}}=V_R.
    \]
\end{proof}

\subsection{The local algebraic Eilenberg theorem}
\label{subsec:to-local-algebraic-eilenberg-theorem}

\begin{theorem}[Local algebraic Eilenberg theorem]
\label{thm:to-local-algebraic-eilenberg}
    The assignments
    \[
        V\longmapsto R_V,
        \qquad
        R\longmapsto V_R
    \]
    restrict to an order-reversing isomorphism between:
    \begin{enumerate}
        \item operation-saturated local coequational theories over
        \((\mathbb A,\mathbb B)\);
        \item operation-closed identity-on-objects operation congruences on
        \(\mathbb A^{\mathrm{op}}\).
    \end{enumerate}
\end{theorem}

\begin{proof}
    If \(V\) is operation-saturated, then
    \[
        V=V_{R_V}.
    \]
    Applying \(R_{(-)}\) and using
    Lemma~\ref{lem:to-operation-closure-operators} gives
    \[
        R_V=R_{V_{R_V}}.
    \]
    Hence \(R_V\) is operation-closed.

    If \(R\) is operation-closed, then
    \[
        R=R_{V_R}.
    \]
    Applying \(V_{(-)}\) and using
    Lemma~\ref{lem:to-operation-closure-operators} gives
    \[
        V_R=V_{R_{V_R}}.
    \]
    Hence \(V_R\) is operation-saturated.

    On fixed points, the composites are identities:
    \[
        V\mapsto R_V\mapsto V_{R_V}=V,
    \]
    and
    \[
        R\mapsto V_R\mapsto R_{V_R}=R.
    \]
    Order reversal follows from
    Proposition~\ref{prop:to-operation-image-contravariant} and
    Lemma~\ref{lem:to-VR-order-reversing}.
\end{proof}

\section{Global algebraic Eilenberg duality}
\label{sec:global-algebraic-eilenberg-duality}

\subsection{Reindexing local theories}
\label{subsec:to-reindexing-local-theories}

We now allow the input category to vary in a chosen category
\[
    \mathsf{Inp}
\]
of internal input categories and internal functors.

Let
\[
    F:\mathbb A'\to\mathbb A
\]
be an internal functor. Recall that residual behaviour reindexing is the map
\[
    F^\sharp:
    A'_0\times_{F_0,A_0,p_{\mathbf{Beh}}}
    \mathbf{Beh}_{\mathbb A,\mathbb B,0}
    \to
    \mathbf{Beh}_{\mathbb A',\mathbb B,0}
\]
given by
\[
    F^\sharp(v',H)=H\circ F_{/v'}.
\]
The domain is regarded over
\[
    Z_{\mathbb A'}=A'_0\times B_0
\]
by
\[
    (v',H)\longmapsto(v',q_{\mathbf{Beh}}H).
\]

\begin{definition}[Reindexed local theory]
\label{def:to-reindexed-local-theory}
    For a local coequational theory
    \[
        V\hookrightarrow
        \mathbf{Beh}^{\mathrm{can}}_{\mathbb A,\mathbb B},
    \]
    define
    \[
        F^\bullet V
        \hookrightarrow
        \mathbf{Beh}^{\mathrm{can}}_{\mathbb A',\mathbb B}
    \]
    to be the regular image of the restricted reindexing map
    \[
        A'_0\times_{F_0,A_0,p_V}V
        \longrightarrow
        \mathbf{Beh}_{\mathbb A',\mathbb B,0}.
    \]
    The domain is regarded over
    \[
        Z_{\mathbb A'}=A'_0\times B_0
    \]
    by
    \[
        (v',H)\longmapsto(v',q_VH).
    \]
\end{definition}

The reindexing image is displayed by
\[
\begin{tikzcd}[column sep=large, row sep=large]
    A'_0\times_{A_0}V
        \arrow[rr, "{F^\sharp|_V}"]
        \arrow[dr, two heads]
    &&
    \mathbf{Beh}_{\mathbb A',\mathbb B,0}
    \\
    &
    F^\bullet V
        \arrow[ur, hook]
    &
\end{tikzcd}
\qquad
\text{in }\mathcal E/Z_{\mathbb A'}.
\]

\begin{lemma}[Reindexed theories are local]
\label{lem:to-reindexed-theories-local}
    For every local coequational theory \(V\) over \((\mathbb A,\mathbb B)\),
    the subobject \(F^\bullet V\) is a local coequational theory over
    \((\mathbb A',\mathbb B)\).
\end{lemma}

\begin{proof}
    The object
    \[
        A'_0\times_{F_0,A_0,p_V}V
    \]
    is the carrier of the input pullback Mealy morphism
    \[
        F^*V:\mathbb A'\rightsquigarrow\mathbb B.
    \]
    Since \(V\) is a restricted canonical object, its semantic map is its
    inclusion into
    \[
        \mathbf{Beh}^{\mathrm{can}}_{\mathbb A,\mathbb B}.
    \]
    The input-pullback semantics theorem gives
    \[
        \beta_{F^*V}(v',H)=F^\sharp(v',H).
    \]
    Therefore the restricted reindexing map is the semantic map of \(F^*V\).
    Its regular image is a local coequational theory by the semantic-image
    theorem.
\end{proof}

\begin{lemma}[Direct images preserve joins]
\label{lem:to-direct-images-preserve-joins}
    In a topos slice, for every morphism
    \[
        f:X\to Y,
    \]
    the image operation
    \[
        \exists_f:\operatorname{Sub}(X)\to\operatorname{Sub}(Y)
    \]
    preserves small joins.
\end{lemma}

\begin{proof}
    The image operation
    \[
        \exists_f
    \]
    is left adjoint to pullback
    \[
        f^*
    \]
    on subobject lattices. Hence it preserves all existing joins.
\end{proof}

\begin{lemma}[Basic properties of \(F^\bullet\)]
\label{lem:to-basic-properties-Fbullet}
    The operation
    \[
        V\longmapsto F^\bullet V
    \]
    is monotone and preserves small joins:
    \[
        V\leq W
        \Longrightarrow
        F^\bullet V\leq F^\bullet W,
    \]
    and
    \[
        F^\bullet\left(\bigvee_i V_i\right)
        =
        \bigvee_i F^\bullet V_i.
    \]
\end{lemma}

\begin{proof}
    If
    \[
        V\leq W,
    \]
    then the restricted reindexing map for \(V\) factors through the restricted
    reindexing map for \(W\). Regular images are monotone on subobjects, so
    \[
        F^\bullet V\leq F^\bullet W.
    \]

    Pullback along
    \[
        A'_0\to A_0
    \]
    preserves joins in the relevant topos slice. The direct image along
    \[
        F^\sharp
    \]
    preserves joins by Lemma~\ref{lem:to-direct-images-preserve-joins}.
    Therefore
    \[
        F^\bullet\left(\bigvee_i V_i\right)
        =
        \bigvee_i F^\bullet V_i.
    \]
\end{proof}

\subsection{Operational inverse image of congruences}
\label{subsec:to-operational-inverse-image}

\begin{definition}[Operational inverse image of a congruence]
\label{def:to-operational-inverse-image-congruence}
    Let
    \[
        R\in
        \mathsf{Cong}_{\mathrm{ioo}}(\mathbb A^{\mathrm{op}}).
    \]
    Its \textbf{operational inverse image} along
    \[
        F:\mathbb A'\to\mathbb A
    \]
    is the operation congruence
    \[
        F^\flat R
        :=
        R_{F^\bullet(V_R)}
    \]
    on
    \[
        (\mathbb A')^{\mathrm{op}}.
    \]
\end{definition}

This is not, in general, the literal pullback relation on arrows. It is the
operation congruence induced after passing through the local Galois closure on
coequational theories.

\begin{lemma}[Monotonicity of operational inverse image]
\label{lem:to-Fflat-monotone}
    If
    \[
        R\leq S
    \]
    in
    \[
        \mathsf{Cong}_{\mathrm{ioo}}(\mathbb A^{\mathrm{op}}),
    \]
    then
    \[
        F^\flat R\leq F^\flat S
    \]
    in
    \[
        \mathsf{Cong}_{\mathrm{ioo}}((\mathbb A')^{\mathrm{op}}).
    \]
\end{lemma}

\begin{proof}
    If
    \[
        R\leq S,
    \]
    then Lemma~\ref{lem:to-VR-order-reversing} gives
    \[
        V_S\leq V_R.
    \]
    Applying monotonicity of \(F^\bullet\),
    \[
        F^\bullet V_S\leq F^\bullet V_R.
    \]
    Applying contravariance of \(R_{(-)}\),
    \[
        R_{F^\bullet V_R}
        \leq
        R_{F^\bullet V_S}.
    \]
    Hence
    \[
        F^\flat R\leq F^\flat S.
    \]
\end{proof}

\begin{lemma}[Raw pullback is contained in operational inverse image]
\label{lem:to-raw-pullback-contained-Fflat}
    Let
    \[
        F:\mathbb A'\to\mathbb A
    \]
    be an input functor, and let \(R\) be an operation congruence on
    \(\mathbb A^{\mathrm{op}}\). Define the raw pullback relation as the pullback of \(R\) along the arrow map
    of
    \[
        F^{\mathrm{op}}:(\mathbb A')^{\mathrm{op}}\to\mathbb A^{\mathrm{op}}.
    \]
    In generalized elements, this says
    \[
        a'\,(F^{-1}R)\,b'
        \quad\Longleftrightarrow\quad
        F_1(a')\,R\,F_1(b')
    \]
    for parallel arrows \(a',b'\) of \(\mathbb A'\). Then \(F^{-1}R\) is an
    identity-on-objects operation congruence on \((\mathbb A')^{\mathrm{op}}\),
    and
    \[
        F^{-1}R\leq F^\flat R.
    \]
\end{lemma}

\begin{proof}
    Since \(R\) is an effective equivalence relation in the arrow slice and
    effective equivalence relations are stable under pullback in a topos, the
    raw pullback is again an effective equivalence relation over
    \[
        A'_0\times A'_0.
    \]
    Compatibility with identities and composition follows from functoriality of
    \(F\) and compatibility of \(R\).

    Let
    \[
        a',b':u'\to v'
    \]
    be parallel arrows with
    \[
        F_1(a')\,R\,F_1(b').
    \]
    Since \(V_R\) is recognized by \(R\), the arrows
    \[
        F_1(a')
        \quad\text{and}\quad
        F_1(b')
    \]
    induce the same transition--output operation on \(V_R\).

    The local theory \(F^\bullet(V_R)\) is the regular image of
    \[
        A'_0\times_{A_0}V_R
        \to
        \mathbf{Beh}_{\mathbb A',\mathbb B,0}.
    \]
    On this regular-epimorphic cover, the operations induced by \(a'\) and
    \(b'\) are the reindexed operations induced by \(F_1(a')\) and
    \(F_1(b')\). Hence they agree on the cover.

    By Lemma~\ref{lem:to-regular-epi-detects-kleisli-equality}, equality
    descends to
    \[
        F^\bullet(V_R).
    \]
    Therefore
    \[
        a'\,R_{F^\bullet(V_R)}\,b',
    \]
    i.e.
    \[
        F^{-1}R\leq F^\flat R.
    \]
\end{proof}

\begin{remark}[Raw pullback versus operational inverse image]
\label{rem:to-raw-pullback-versus-operational-pullback}
    The congruence \(F^\flat R\) is the operation congruence generated by
    residual reindexing of the largest theory recognized by \(R\). The raw
    pullback congruence \(F^{-1}R\) is contained in it by
    Lemma~\ref{lem:to-raw-pullback-contained-Fflat}. Equality holds when
    residual behaviour reindexing reflects the relevant transition--output
    distinctions.
\end{remark}

\subsection{Global theories and global congruences}
\label{subsec:to-global-theories-congruences}

\begin{definition}[Global operation congruence theory]
\label{def:to-global-operation-congruence-theory}
    A \textbf{global operation congruence theory} consists of an operation
    congruence
    \[
        R_{\mathbb A}\in
        \mathsf{Cong}_{\mathrm{ioo}}(\mathbb A^{\mathrm{op}})
    \]
    for each input category \(\mathbb A\in\mathsf{Inp}\), such that for every
    internal functor
    \[
        F:\mathbb A'\to\mathbb A
    \]
    one has
    \[
        R_{\mathbb A'}\leq F^\flat R_{\mathbb A}.
    \]
\end{definition}

Global coequational theories are ordered pointwise by inclusion. Global
operation congruence theories are ordered pointwise by inclusion of
congruences.

\begin{definition}[Global saturation and closure]
\label{def:to-global-operation-saturation-closure}
    A global coequational theory
    \[
        V=(V_{\mathbb A})_{\mathbb A}
    \]
    is \textbf{operation-saturated} if each \(V_{\mathbb A}\) is
    operation-saturated:
    \[
        V_{\mathbb A}=V_{R_{V_{\mathbb A}}}.
    \]
    A global operation congruence theory
    \[
        R=(R_{\mathbb A})_{\mathbb A}
    \]
    is \textbf{operation-closed} if each \(R_{\mathbb A}\) is operation-closed:
    \[
        R_{\mathbb A}=R_{V_{R_{\mathbb A}}}.
    \]
\end{definition}

\begin{proposition}[Saturated global theories induce closed global congruences]
\label{prop:to-saturated-global-theories-to-closed-global-congruences}
    If
    \[
        V=(V_{\mathbb A})_{\mathbb A\in\mathsf{Inp}}
    \]
    is an operation-saturated global coequational theory, then
    \[
        R^V_{\mathbb A}:=R_{V_{\mathbb A}}
    \]
    defines an operation-closed global operation congruence theory.
\end{proposition}

\begin{proof}
    Local operation-closure follows from
    Theorem~\ref{thm:to-local-algebraic-eilenberg}. Let
    \[
        F:\mathbb A'\to\mathbb A
    \]
    be an internal functor. Since \(V\) is global,
    \[
        F^\bullet V_{\mathbb A}\leq V_{\mathbb A'}.
    \]
    Contravariance gives
    \[
        R_{V_{\mathbb A'}}
        \leq
        R_{F^\bullet V_{\mathbb A}}.
    \]
    Since \(V_{\mathbb A}\) is operation-saturated,
    \[
        V_{\mathbb A}=V_{R_{V_{\mathbb A}}}.
    \]
    Therefore
    \[
        R_{F^\bullet V_{\mathbb A}}
        =
        R_{F^\bullet(V_{R_{V_{\mathbb A}}})}
        =
        F^\flat R_{V_{\mathbb A}}.
    \]
    Hence
    \[
        R^V_{\mathbb A'}\leq F^\flat R^V_{\mathbb A}.
    \]
\end{proof}

\begin{proposition}[Closed global congruences induce saturated global theories]
\label{prop:to-closed-global-congruences-to-saturated-global-theories}
    If
    \[
        R=(R_{\mathbb A})_{\mathbb A\in\mathsf{Inp}}
    \]
    is an operation-closed global operation congruence theory, then
    \[
        V^R_{\mathbb A}:=V_{R_{\mathbb A}}
    \]
    defines an operation-saturated global coequational theory.
\end{proposition}

\begin{proof}
    Local operation-saturation follows from
    Theorem~\ref{thm:to-local-algebraic-eilenberg}. Let
    \[
        F:\mathbb A'\to\mathbb A
    \]
    be an internal functor. We prove
    \[
        F^\bullet V^R_{\mathbb A}\leq V^R_{\mathbb A'}.
    \]
    The global congruence condition gives
    \[
        R_{\mathbb A'}\leq F^\flat R_{\mathbb A}
        =
        R_{F^\bullet(V_{R_{\mathbb A}})}.
    \]
    By Theorem~\ref{thm:to-local-operation-galois}, this is equivalent to
    \[
        F^\bullet(V_{R_{\mathbb A}})
        \leq
        V_{R_{\mathbb A'}}.
    \]
    Hence
    \[
        F^\bullet V^R_{\mathbb A}\leq V^R_{\mathbb A'}.
    \]
    Thus \(V^R\) is global.
\end{proof}

\subsection{The global algebraic Eilenberg theorem}
\label{subsec:to-global-algebraic-eilenberg-theorem}

\begin{theorem}[Global algebraic Eilenberg theorem]
\label{thm:to-global-algebraic-eilenberg}
    The assignments
    \[
        V=(V_{\mathbb A})_{\mathbb A}
        \longmapsto
        R^V=(R_{V_{\mathbb A}})_{\mathbb A}
    \]
    and
    \[
        R=(R_{\mathbb A})_{\mathbb A}
        \longmapsto
        V^R=(V_{R_{\mathbb A}})_{\mathbb A}
    \]
    restrict to an order-reversing isomorphism between:
    \begin{enumerate}
        \item operation-saturated global coequational theories;
        \item operation-closed global operation congruence theories.
    \end{enumerate}
\end{theorem}

\begin{proof}
    Proposition~\ref{prop:to-saturated-global-theories-to-closed-global-congruences}
    sends saturated global coequational theories to closed global congruence
    theories. Proposition~\ref{prop:to-closed-global-congruences-to-saturated-global-theories}
    sends closed global congruence theories to saturated global coequational
    theories. The fact that the assignments are mutually inverse on fixed
    points is the local fixed-point theorem applied pointwise. Order reversal
    is also pointwise, using the local order-reversing correspondence.
\end{proof}

\section{The one-object specialization}
\label{sec:to-one-object-specialization}

This section records the one-object form of the transition--output image
construction. Let the input category \(\mathbb A\) be determined by a monoid
object \(M\). We keep the standing convention
\[
    b\circ a=ba .
\]
Thus a right action law appears as
\[
    x\cdot(mn)=(x\cdot m)\cdot n.
\]
Passing to
\[
    \mathbb A^{\mathrm{op}}
\]
restores the usual multiplication order on the operation side.

\subsection{The one-object Kleisli endomorphism monoid}
\label{subsec:to-one-object-klend}

In the one-object case, a carrier over the unique input object is simply an
object
\[
    q:P\to B_0.
\]
The internal category
\[
    \operatorname{KlEnd}_{\mathbb B}(P)
\]
has one object, and its arrow object is the internal monoid of
output-comparison Kleisli endomorphisms
\[
    P\rightsquigarrow P.
\]
A Kleisli endomorphism consists of an endomap
\[
    F:P\to P
\]
together with an output-comparison arrow
\[
    \alpha_x:q(Fx)\to q(x)
\]
in \(\mathbb B\).

If \(\mathbb B\) is also one-object, determined by a monoid object \(N\), then
a Kleisli endomorphism is an endomap of \(P\) together with an \(N\)-valued
output-comparison map. The Kleisli composite of
\[
    (F,\alpha),\qquad (G,\gamma)
\]
is
\[
    (G,\gamma)\odot(F,\alpha)
    =
    (GF,\alpha\circ_F\gamma),
\]
where
\[
    (\alpha\circ_F\gamma)_x
    =
    \alpha_x\gamma_{Fx}.
\]
Thus the output component is multiplied in the order dictated by the
comparison chain
\[
    q(GFx)\xrightarrow{\gamma_{Fx}}q(Fx)
    \xrightarrow{\alpha_x}q(x).
\]

\subsection{One-object Mealy actions}
\label{subsec:to-one-object-mealy-actions}

A one-object Mealy morphism is a right \(M\)-object \(P\), written
\[
    x\cdot m:=\delta(m,x),
\]
together with an output cocycle
\[
    \omega:M\times P\to N
\]
satisfying
\[
    \omega(1,x)=1,
    \qquad
    \omega(mn,x)=\omega(m,x)\omega(n,x\cdot m).
\]
The transition--output action is the monoid homomorphism
\[
    \chi_P:
    M^{\mathrm{op}}\to
    \operatorname{KlEnd}_{\mathbb B}(P)
\]
sending
\[
    m\longmapsto
    \bigl(x\mapsto x\cdot m,\;\omega(m,x)\bigr).
\]
The appearance of \(M^{\mathrm{op}}\) compensates for the target-to-source
processing convention.

\subsection{Residual behaviours and local theories}
\label{subsec:to-one-object-residual-theories}

The residual category of \(M\) is the right Cayley residual category. Its
objects are elements \(r\in M\), and a residual arrow labelled \(m\) has the
form
\[
    rm\longrightarrow r .
\]
When the output category is one-object with output monoid \(N\), a residual
behaviour is an \(N\)-valued cocycle
\[
    \lambda:M\times M\to N
\]
satisfying
\[
    \lambda(r,1)=1,
    \qquad
    \lambda(r,mn)=\lambda(r,m)\lambda(rm,n).
\]
The canonical derivative is left translation in the residual-history variable:
\[
    (\partial_s\lambda)(r,m)=\lambda(sr,m),
\]
and the canonical output on input \(s\) is
\[
    \lambda(1,s).
\]

A local coequational theory in the one-object case is a subobject
\[
    V\hookrightarrow \mathbf{Beh}_0
\]
of residual behaviours closed under all derivatives:
\[
    \lambda\in V
    \quad\Longrightarrow\quad
    \partial_s\lambda\in V
    \qquad(s\in M).
\]
The residual transition--output action of \(V\) is
\[
    \chi_V:
    M^{\mathrm{op}}\to
    \operatorname{KlEnd}_{\mathbb B}(V),
\]
with
\[
    \chi_V(s)(\lambda)
    =
    \bigl(\partial_s\lambda,\lambda(1,s)\bigr).
\]

\subsection{One-object transition--output images}
\label{subsec:to-one-object-operation-images}

For a local coequational theory
\[
    V\hookrightarrow\mathbf{Beh}^{\mathrm{can}},
\]
the transition--output image
\[
    \operatorname{Op}(V)
\]
is a quotient monoid of \(M^{\mathrm{op}}\). Its congruence identifies two
elements \(m,n\in M\) precisely when they induce the same transition--output
Kleisli endomorphism on all states of \(V\). Thus
\[
    m\,R_V\,n
\]
means that for every \(\lambda\in V\),
\[
    \partial_m\lambda=\partial_n\lambda
\]
and the corresponding output comparisons agree:
\[
    \lambda(1,m)=\lambda(1,n).
\]
In cases where the output category is chaotic, the output-comparison condition
is determined by the endpoint values and hence carries no additional arrow
data.

For a specification
\[
    k:K\to\mathbf{Beh}_0,
\]
the syntactic transition--output monoid is
\[
    \operatorname{SynTO}_{M,\mathbb B}(K)
    :=
    \operatorname{Op}(\operatorname{Der}_0(K)).
\]
Equivalently, it is the quotient of \(M^{\mathrm{op}}\) by the congruence
that identifies two elements exactly when they have the same transition--output
effect on all residual behaviours generated by \(K\).

\subsection{One-object algebraic Eilenberg duality}
\label{subsec:to-one-object-eilenberg}

In the one-object setting, the local operation Galois theorem specializes to
an order-reversing correspondence between derivative-closed behaviour theories
and operation congruences on \(M^{\mathrm{op}}\), up to the fixed-point
conditions:
\[
    V=V_{R_V},
    \qquad
    R=R_{V_R}.
\]
Thus a derivative-closed local theory determines the congruence of operations
it cannot distinguish, and a congruence determines the largest derivative-
closed theory whose transition--output action factors through the quotient.

When the input monoid varies, a monoid homomorphism
\[
    h:M'\to M
\]
acts on residual behaviours by inverse reindexing. In the free-monoid Boolean
case, this becomes ordinary inverse morphic preimage of languages:
\[
    h^{-1}L=L\circ h.
\]
The global operation congruence condition says that the congruence assigned to
\(M'\) is contained in the operational inverse image of the congruence assigned
to \(M\).

\section{The stateless specialization}
\label{sec:to-stateless-specialization}

This section records the stateless form of the transition--output image
construction. A stateless Mealy morphism is induced by an internal functor
\[
    F:\mathbb A\to\mathbb B.
\]
Its carrier is
\[
    A_0\xleftarrow{1_{A_0}}A_0\xrightarrow{F_0}B_0.
\]
Thus the fibre over an input object \(v\in A_0\) is terminal.

\subsection{The stateless Kleisli endomorphism category}
\label{subsec:to-stateless-klend}

For the stateless carrier
\[
    P=A_0,
    \qquad
    q=F_0,
\]
a Kleisli arrow
\[
    P_v\rightsquigarrow P_u
\]
is determined by a single arrow
\[
    F_0(u)\to F_0(v)
\]
of \(\mathbb B\). Hence
\[
    \operatorname{KlEnd}_{\mathbb B}(A_0)
\]
records possible output-comparison arrows between the object values of \(F\).

The transition--output action sends an arrow
\[
    a:u\to v
\]
of \(\mathbb A\) to the unique state transition
\[
    v\longmapsto u
\]
together with the output comparison
\[
    F_1(a):F_0(u)\to F_0(v).
\]
Thus the stateless transition--output action is exactly the functor \(F\),
viewed through the opposite input orientation:
\[
    \chi_F:
    \mathbb A^{\mathrm{op}}
    \to
    \operatorname{KlEnd}_{\mathbb B}(A_0).
\]

\subsection{Stateless residual behaviours}
\label{subsec:to-stateless-residual-behaviours}

The residual behaviour of the stateless functor \(F\) at an object
\[
    v\in A_0
\]
is the residual slice functor
\[
    H^F_v:\mathbb A/v\to\mathbb B.
\]
On a residual object
\[
    r:u\to v,
\]
it is given by
\[
    H^F_v(r)=F_0(u),
\]
and on a residual morphism represented by a triangle
\[
\begin{tikzcd}[column sep=large, row sep=large]
    w
        \arrow[r, "a"]
        \arrow[rr, "r\circ a", bend left=15]
    &
    u
        \arrow[r, "r"']
    &
    v
\end{tikzcd}
\]
it is given by
\[
    H^F_v(r\circ a\to r)=F_1(a).
\]
Thus residual semantics of a stateless functor records all slice-restricted
views of \(F\).

The derivative along
\[
    a:u\to v
\]
sends
\[
    H^F_v
\]
to
\[
    H^F_u,
\]
and the canonical output comparison is
\[
    F_1(a):F_0(u)\to F_0(v).
\]

\subsection{Stateless transition--output images}
\label{subsec:to-stateless-operation-images}

For a stateless functor
\[
    F:\mathbb A\to\mathbb B,
\]
the transition--output image
\[
    \operatorname{Op}(F)
\]
is the identity-on-objects regular image of
\[
    \mathbb A^{\mathrm{op}}
    \to
    \operatorname{KlEnd}_{\mathbb B}(A_0)
\]
induced by \(F\). It identifies parallel arrows
\[
    a,a':u\to v
\]
whenever they induce the same output comparison:
\[
    F_1(a)=F_1(a').
\]
Thus, in the pure stateless case, the transition--output image is the
identity-on-objects image of \(F\) on arrows, with the input orientation
reversed.

More generally, if a stateless specification is generated by a family of
internal functors
\[
    F_i:\mathbb A\to\mathbb B,
\]
then the generated local coequational theory contains the residual slice
functors
\[
    H^{F_i}_v:\mathbb A/v\to\mathbb B.
\]
The associated transition--output congruence identifies parallel arrows
\[
    a,a':u\to v
\]
precisely when all generated stateless residual behaviours assign the same
output comparison to them. For the immediate stateless states this includes
the equations
\[
    F_{i,1}(a)=F_{i,1}(a')
\]
for all \(i\).

Residual contexts determine which slice-restricted states occur in the
generated theory. In the stateless case, they do not add extra output data
beyond the functorial arrow images themselves; rather, they organize those
arrow images as residual behaviours over the slices \(\mathbb A/v\).

\subsection{Stateless syntactic transition--output categories}
\label{subsec:to-stateless-syntactic-categories}

For a stateless specification generated by functors
\[
    F_i:\mathbb A\to\mathbb B,
\]
the syntactic transition--output category is obtained by quotienting
\[
    \mathbb A^{\mathrm{op}}
\]
by equality of the transition--output operations induced on all generated
residual slice functors. It is therefore the coarsest identity-on-objects
quotient of \(\mathbb A^{\mathrm{op}}\) through which the specified residual
output functors factor.

In the case of a single functor
\[
    F:\mathbb A\to\mathbb B,
\]
this quotient identifies input arrows exactly when they induce the same output
arrow under \(F\). Thus the stateless syntactic transition--output category is
the image of the functorial output action, expressed in the same
transition--output language as the general Mealy case.

\subsection{Stateless globality}
\label{subsec:to-stateless-globality}

In the stateless setting, input reindexing is precomposition. If
\[
    G:\mathbb A\to\mathbb B
\]
is a stateless system and
\[
    F:\mathbb A'\to\mathbb A
\]
is an input functor, then the input pullback is the stateless system
\[
    G\circ F:\mathbb A'\to\mathbb B.
\]
Hence globality says that the associated operation congruences are stable under
precomposition of input categories.

Thus if
\[
    G:\mathbb A\to\mathbb B
\]
satisfies a saturated global theory, then so does every composite
\[
    \mathbb A'\xrightarrow{F}\mathbb A\xrightarrow{G}\mathbb B.
\]
On the congruence side, this is expressed by the operational inverse-image
condition
\[
    R_{\mathbb A'}\leq F^\flat R_{\mathbb A}.
\]

\section{Examples}
\label{sec:to-examples}

This section collects the main concrete examples of the transition--output
image construction.

\subsection{Boolean output}
\label{subsec:to-boolean-output-example}

For Boolean output, take
\[
    \mathbb B=2_{\mathrm{ch}},
\]
the chaotic internal category on a Boolean algebra object \(2\). The Boolean
operations
\[
    \wedge,\vee,\neg,0,1
\]
are internal functors
\[
    2_{\mathrm{ch}}^n\to 2_{\mathrm{ch}}.
\]
Since \(2_{\mathrm{ch}}\) is chaotic, an output-comparison arrow is uniquely
determined by its endpoints.

In \(\mathbf{Set}\), a residual behaviour for the free-monoid input category
is determined by its object assignment, hence by a language
\[
    L:\Sigma^*\to 2.
\]
The derivative convention is
\[
    (\partial_uL)(w)=L(uw).
\]

A plain local Boolean coequational theory is a derivative-closed collection of
languages. If the Boolean algebraic structure is imposed, then a local Boolean
coequational theory is a Boolean algebra of languages closed under residual
derivatives. It becomes a Boolean algebra of regular languages after imposing a
finite-state condition or after restricting the behaviour object to regular
behaviours. The derivative notation is the one introduced by Brzozowski for
regular expressions and languages
\cite{Brzozowski1964}.

\subsection{The one-object free-monoid case}
\label{subsec:to-free-monoid-specialization}

Let
\[
    \mathbb A_\Sigma
\]
be the one-object internal category associated to the free monoid
\[
    \Sigma^*.
\]
With the one-object convention of
Chapter~\ref{ch:spans-internal-categories-mealy}, categorical composition is
written
\[
    b\circ a=ba.
\]
Thus
\[
    (\partial_uL)(w)=L(uw).
\]

Passing to
\[
    \mathbb A_\Sigma^{\mathrm{op}}
\]
restores the ordinary word order. The one-object category
\(\mathbb A_\Sigma\) has categorical multiplication
\[
    m_{\mathbb A_\Sigma}(u,v)=vu,
\]
while the arrow monoid of
\[
    \mathbb A_\Sigma^{\mathrm{op}}
\]
has multiplication
\[
    m_{\mathbb A_\Sigma^{\mathrm{op}}}(u,v)=uv.
\]

\begin{proposition}[Classical syntactic monoid recovery]
\label{prop:to-classical-syntactic-monoid-recovery}
    Let
    \[
        L\subseteq\Sigma^*
    \]
    be a regular language, regarded as a specification
    \[
        1\to
        \mathbf{Beh}_{\mathbb A_\Sigma,2_{\mathrm{ch}},0}.
    \]
    Then the arrow monoid of
    \[
        \operatorname{SynTO}_{\mathbb A_\Sigma,2_{\mathrm{ch}}}(L)
        =
        \operatorname{Op}(\operatorname{Der}_0(L))
    \]
    is canonically isomorphic to the classical syntactic monoid of \(L\).
\end{proposition}

\begin{proof}
    In \(2_{\mathrm{ch}}\), there is a unique arrow between any two Boolean
    values. Therefore, once the transformed residual languages agree, the
    output-arrow equality in Theorem~\ref{thm:to-context-expansion} is
    automatic. The congruence calculation below applies to any language;
    regularity is used to ensure that the resulting syntactic monoid is finite.

    The residual state
    \[
        \partial_rL
    \]
    represents the left quotient by \(r\). Acting by \(u\) sends it to
    \[
        \partial_{ru}L.
    \]
    Equality of the operations induced by \(u\) and \(v\) on all residual
    states therefore says
    \[
        \partial_{ru}L=\partial_{rv}L
        \qquad(r\in\Sigma^*).
    \]
    Evaluating both sides at an arbitrary continuation \(s\in\Sigma^*\) gives
    \[
        L(rus)=L(rvs).
    \]
    Thus
    \[
        u\equiv v
        \quad\Longleftrightarrow\quad
        \forall r,s\in\Sigma^*,\;
        rus\in L\Longleftrightarrow rvs\in L.
    \]
    This is the classical syntactic congruence of \(L\).

    The arrow monoid of
    \[
        \operatorname{SynTO}_{\mathbb A_\Sigma,2_{\mathrm{ch}}}(L)
    \]
    is the quotient of the arrow monoid of
    \[
        \mathbb A_\Sigma^{\mathrm{op}}
    \]
    by this congruence. Since the passage to
    \[
        \mathbb A_\Sigma^{\mathrm{op}}
    \]
    cancels the convention
    \[
        b\circ a=ba,
    \]
    the quotient multiplication is ordinary concatenation. Hence the arrow
    monoid is canonically the classical syntactic monoid of \(L\), as in
    Eilenberg's algebraic treatment of regular languages
    \cite{Eilenberg1974}.
\end{proof}

\subsection{Compatibility with the classical Eilenberg correspondence}
\label{subsec:to-classical-eilenberg-compatibility}

The preceding proposition identifies the transition--output image of an
individual regular language with its classical syntactic monoid. Passing from
individual languages to Eilenberg varieties uses Boolean closure, two-sided
quotient closure, inverse morphic preimage closure, and finite-state
recognizability
\cite{Eilenberg1974}.

\begin{definition}[Global Boolean regular language theory]
\label{def:to-global-boolean-regular-language-theory}
    A \textbf{global Boolean regular language theory} is a family
    \[
        \mathcal V=(\mathcal V_\Sigma)_\Sigma
    \]
    such that:
    \begin{enumerate}
        \item each \(\mathcal V_\Sigma\) is a Boolean algebra of regular
        languages over the finite alphabet \(\Sigma\);
        \item each \(\mathcal V_\Sigma\) is closed under left and right
        quotients;
        \item for every monoid homomorphism
        \[
            h:\Gamma^*\to\Sigma^*,
        \]
        inverse image sends
        \[
            \mathcal V_\Sigma
            \quad\text{into}\quad
            \mathcal V_\Gamma.
        \]
    \end{enumerate}
\end{definition}

Given such a theory \(\mathcal V\), define \(\mathsf P_{\mathcal V}\) to be
the class of finite monoids dividing finite products of syntactic monoids
\[
    \operatorname{Syn}(L),
    \qquad
    L\in\mathcal V_\Sigma,
\]
as \(\Sigma\) varies. Conversely, for a pseudovariety of finite monoids
\(\mathsf P\), define
\[
    \mathcal V^{\mathsf P}_\Sigma
    =
    \{L\subseteq\Sigma^*
      \mid
      \operatorname{Syn}(L)\in\mathsf P\}.
\]

\begin{lemma}[Syntactic separation]
\label{lem:to-syntactic-separation}
    Let \(L\subseteq\Sigma^*\) be regular. Every fibre of the syntactic
    morphism of \(L\) is a Boolean combination of two-sided quotients
    \[
        x^{-1}Ly^{-1}
        =
        \{w\in\Sigma^*\mid xwy\in L\}.
    \]
\end{lemma}

\begin{proof}
    Two words lie in the same syntactic congruence class precisely when no
    two-sided context
    \[
        x(-)y
    \]
    separates them with respect to membership in \(L\). Since the syntactic
    monoid is finite, each congruence class is the finite intersection of the
    two-sided quotient conditions true on that class and the complements of
    those that fail on that class. Hence each fibre is a Boolean combination of
    two-sided quotients of \(L\).
\end{proof}

\begin{lemma}[Finite syntactic generation]
\label{lem:to-finite-syntactic-generation}
    Let \(M\) be a finite monoid generated by a finite alphabet \(\Sigma\).
    Then \(M\) divides a finite product of syntactic monoids of languages
    recognized by \(M\).
\end{lemma}

\begin{proof}
    Choose a surjective monoid morphism
    \[
        h:\Sigma^*\twoheadrightarrow M.
    \]
    For each subset \(S\subseteq M\), the language
    \[
        h^{-1}(S)
    \]
    is recognized by \(M\), and its syntactic monoid divides \(M\). Choose
    finitely many subsets of \(M\) separating distinct elements of \(M\). The
    product of the corresponding syntactic morphisms separates all elements of
    \(M\). Hence \(M\) embeds into a finite product of these syntactic monoids,
    and therefore \(M\) divides that finite product.
\end{proof}

\begin{lemma}[Recognition by products of syntactic monoids]
\label{lem:to-products-of-syntactic-monoids-recognize-variety-languages}
    Let \(\mathcal V\) be a global Boolean regular language theory. If a
    language
    \[
        L\subseteq\Sigma^*
    \]
    is recognized by a divisor of a finite product of syntactic monoids
    \[
        \operatorname{Syn}(L_i),
        \qquad
        L_i\in\mathcal V_{\Sigma_i},
    \]
    then
    \[
        L\in\mathcal V_\Sigma.
    \]
\end{lemma}

\begin{proof}
    A divisor is a quotient of a submonoid of a finite product. Since
    \(\Sigma^*\) is free, each component morphism
    \[
        \Sigma^*\to\operatorname{Syn}(L_i)
    \]
    lifts through the syntactic morphism
    \[
        \eta_{L_i}:\Sigma_i^*\twoheadrightarrow\operatorname{Syn}(L_i)
    \]
    after choosing, for every generator \(\sigma\in\Sigma\), a word in
    \(\Sigma_i^*\) mapping to the required component value. Thus each
    component is induced by a monoid homomorphism
    \[
        h_i:\Sigma^*\to\Sigma_i^*.
    \]

    By Lemma~\ref{lem:to-syntactic-separation}, fibres of each
    \(\eta_{L_i}\) are Boolean combinations of two-sided quotients of \(L_i\).
    Since \(\mathcal V\) is closed under two-sided quotients, Boolean
    operations, and inverse morphic preimages, the inverse images of these
    fibres along the maps \(h_i\) lie in \(\mathcal V_\Sigma\). Finite product
    fibres are finite Boolean combinations of such inverse images, hence also
    lie in \(\mathcal V_\Sigma\).

    A language recognized by a divisor of the product is a union of fibres of
    the induced finite recognizing morphism. Since the monoid is finite, this
    union is finite, hence a Boolean combination of product fibres. Therefore
    \[
        L\in\mathcal V_\Sigma.
    \]
\end{proof}

\begin{theorem}[Compatibility with the classical Eilenberg correspondence]
\label{thm:to-classical-eilenberg-recovery}
    In the global Boolean regular-language specialization, the assignments
    \[
        \mathcal V\longmapsto \mathsf P_{\mathcal V},
        \qquad
        \mathsf P\longmapsto \mathcal V^{\mathsf P}
    \]
    are the classical Eilenberg correspondence between varieties of regular
    languages and pseudovarieties of finite monoids.
\end{theorem}

\begin{proof}
    By Proposition~\ref{prop:to-classical-syntactic-monoid-recovery}, the
    transition--output image associated to a regular language \(L\) has arrow
    monoid equal to the syntactic monoid of \(L\).

    If
    \[
        \mathcal V=(\mathcal V_\Sigma)_\Sigma
    \]
    is a global Boolean regular language theory, then
    \[
        \mathsf P_{\mathcal V}
    \]
    is closed under finite products and divisors by construction. Hence it is a
    pseudovariety.

    If
    \[
        \mathsf P
    \]
    is a pseudovariety, then
    \[
        \mathcal V^{\mathsf P}_\Sigma
        =
        \{L\subseteq\Sigma^*
          \mid
          \operatorname{Syn}(L)\in\mathsf P\}
    \]
    is closed under Boolean operations, quotients, and inverse morphic
    preimages. These operations preserve recognizability by monoids in
    \(\mathsf P\) up to finite products and divisors.

    The inclusion
    \[
        \mathsf P_{\mathcal V^{\mathsf P}}\subseteq\mathsf P
    \]
    follows from closure of \(\mathsf P\) under finite products and divisors.
    The reverse inclusion follows from
    Lemma~\ref{lem:to-finite-syntactic-generation}: every finite monoid in
    \(\mathsf P\) divides a finite product of syntactic monoids of languages
    recognized by monoids in \(\mathsf P\), and those languages belong to
    \(\mathcal V^{\mathsf P}\).

    The inclusion
    \[
        \mathcal V\subseteq\mathcal V^{\mathsf P_{\mathcal V}}
    \]
    is immediate, since every
    \[
        L\in\mathcal V_\Sigma
    \]
    has syntactic monoid among the generators of
    \[
        \mathsf P_{\mathcal V}.
    \]
    Conversely, suppose
    \[
        L\in \mathcal V^{\mathsf P_{\mathcal V}}_\Sigma.
    \]
    Then \(\operatorname{Syn}(L)\) divides a finite product of syntactic
    monoids of languages
    \[
        L_i\in\mathcal V_{\Sigma_i}.
    \]
    By
    Lemma~\ref{lem:to-products-of-syntactic-monoids-recognize-variety-languages},
    this implies
    \[
        L\in\mathcal V_\Sigma.
    \]
    Therefore
    \[
        \mathcal V^{\mathsf P_{\mathcal V}}\subseteq\mathcal V.
    \]
    Hence the two assignments are mutually inverse.
\end{proof}

\subsection{A stateless comparison example}
\label{subsec:to-stateless-comparison-example}

Let
\[
    F,G:\mathbb A\to\mathbb B
\]
be internal functors, regarded as stateless Mealy morphisms. Their generated
local coequational theory contains the residual slice functors
\[
    H^F_v:\mathbb A/v\to\mathbb B,
    \qquad
    H^G_v:\mathbb A/v\to\mathbb B.
\]
For an arrow
\[
    a:u\to v,
\]
the action on the generated stateless residual state \(H^F_v\) is
\[
    H^F_v\longmapsto H^F_u
\]
with output comparison
\[
    F_1(a):F_0(u)\to F_0(v).
\]
Similarly, the action on \(H^G_v\) has output comparison
\[
    G_1(a):G_0(u)\to G_0(v).
\]

Therefore the associated transition--output image identifies parallel arrows
\[
    a,a':u\to v
\]
whenever the generated residual behaviours assign the same transition--output
operation to them. In particular, for the immediate stateless states one tests
\[
    F_1(a)=F_1(a'),
    \qquad
    G_1(a)=G_1(a').
\]
The residual slice description records these tests uniformly over the generated
residual behaviours. Thus the stateless example shows that the
transition--output image is a quotient of input arrows by their action on the
generated residual semantics, and in the pure functorial case this action is
exactly the output-arrow data of the functors themselves.

\section{Summary}
\label{sec:syntactic-to-summary}

This chapter extracted the algebraic side of coequational residual semantics.

\[
\begin{array}{c|l}
\text{construction} & \text{role} \\
\hline
T_{\mathbb B}
&
\text{output-comparison Kleisli monad}
\\
\operatorname{KlEnd}_{\mathbb B}(P)
&
\text{fibrewise Kleisli endomorphism category}
\\
\chi_P
&
\text{residual transition--output action of a Mealy morphism}
\\
\operatorname{Op}(V)
&
\text{transition--output image of a local theory}
\\
R_V
&
\text{operation congruence of }V
\\
V_R
&
\text{largest local theory recognized by }R
\\
\operatorname{SynTO}(K)
&
\text{syntactic transition--output category of a specification}
\\
F^\bullet
&
\text{input reindexing of local theories}
\\
F^\flat
&
\text{operational inverse image of congruences}
\end{array}
\]

A Mealy morphism is equivalently a fibrewise Kleisli action
\[
    \chi_P:
    \mathbb A^{\mathrm{op}}
    \to
    \operatorname{KlEnd}_{\mathbb B}(P).
\]
For a local coequational theory \(V\), the restricted canonical object gives
\[
    \chi_V:
    \mathbb A^{\mathrm{op}}
    \to
    \operatorname{KlEnd}_{\mathbb B}(V).
\]
Its identity-on-objects regular image
\[
    \operatorname{Op}(V)
\]
is the transition--output quotient of the input category.

For a specification \(K\), the syntactic transition--output category is
\[
    \operatorname{SynTO}_{\mathbb A,\mathbb B}(K)
    =
    \operatorname{Op}(\operatorname{Der}_0(K)).
\]
The context expansion theorem identifies its congruence by testing generated
residual contexts.

The local algebraic Eilenberg theorem gives an order-reversing isomorphism
between operation-saturated local coequational theories and operation-closed
identity-on-objects operation congruences. The global theorem extends this to
varying input categories by transporting local theories through
\[
    F^\bullet
\]
and congruences through
\[
    F^\flat R=R_{F^\bullet(V_R)}.
\]

In the one-object Boolean specialization, the arrow monoid of
\[
    \operatorname{SynTO}(L)
\]
is the classical syntactic monoid of \(L\). In the global Boolean regular
language specialization, this recovers the syntactic-monoid side of
Eilenberg's correspondence.